% Options for packages loaded elsewhere
\PassOptionsToPackage{unicode}{hyperref}
\PassOptionsToPackage{hyphens}{url}
%
\documentclass[
  letterpaper,
  paper=7in:10in,
  pagesize=pdftex,
  headinclude=on,
  footinclude=on,
  12pt]{scrbook}

\usepackage{amsmath,amssymb}
\usepackage{iftex}
\ifPDFTeX
  \usepackage[T1]{fontenc}
  \usepackage[utf8]{inputenc}
  \usepackage{textcomp} % provide euro and other symbols
\else % if luatex or xetex
  \usepackage{unicode-math}
  \defaultfontfeatures{Scale=MatchLowercase}
  \defaultfontfeatures[\rmfamily]{Ligatures=TeX,Scale=1}
\fi
\usepackage{lmodern}
\ifPDFTeX\else  
    % xetex/luatex font selection
\fi
% Use upquote if available, for straight quotes in verbatim environments
\IfFileExists{upquote.sty}{\usepackage{upquote}}{}
\IfFileExists{microtype.sty}{% use microtype if available
  \usepackage[]{microtype}
  \UseMicrotypeSet[protrusion]{basicmath} % disable protrusion for tt fonts
}{}
\makeatletter
\@ifundefined{KOMAClassName}{% if non-KOMA class
  \IfFileExists{parskip.sty}{%
    \usepackage{parskip}
  }{% else
    \setlength{\parindent}{0pt}
    \setlength{\parskip}{6pt plus 2pt minus 1pt}}
}{% if KOMA class
  \KOMAoptions{parskip=half}}
\makeatother
\usepackage{xcolor}
\ifLuaTeX
  \usepackage{luacolor}
  \usepackage[soul]{lua-ul}
\else
  \usepackage{soul}
  
\fi
\setlength{\emergencystretch}{3em} % prevent overfull lines
\setcounter{secnumdepth}{5}
% Make \paragraph and \subparagraph free-standing
\makeatletter
\ifx\paragraph\undefined\else
  \let\oldparagraph\paragraph
  \renewcommand{\paragraph}{
    \@ifstar
      \xxxParagraphStar
      \xxxParagraphNoStar
  }
  \newcommand{\xxxParagraphStar}[1]{\oldparagraph*{#1}\mbox{}}
  \newcommand{\xxxParagraphNoStar}[1]{\oldparagraph{#1}\mbox{}}
\fi
\ifx\subparagraph\undefined\else
  \let\oldsubparagraph\subparagraph
  \renewcommand{\subparagraph}{
    \@ifstar
      \xxxSubParagraphStar
      \xxxSubParagraphNoStar
  }
  \newcommand{\xxxSubParagraphStar}[1]{\oldsubparagraph*{#1}\mbox{}}
  \newcommand{\xxxSubParagraphNoStar}[1]{\oldsubparagraph{#1}\mbox{}}
\fi
\makeatother


\providecommand{\tightlist}{%
  \setlength{\itemsep}{0pt}\setlength{\parskip}{0pt}}\usepackage{longtable,booktabs,array}
\usepackage{calc} % for calculating minipage widths
% Correct order of tables after \paragraph or \subparagraph
\usepackage{etoolbox}
\makeatletter
\patchcmd\longtable{\par}{\if@noskipsec\mbox{}\fi\par}{}{}
\makeatother
% Allow footnotes in longtable head/foot
\IfFileExists{footnotehyper.sty}{\usepackage{footnotehyper}}{\usepackage{footnote}}
\makesavenoteenv{longtable}
\usepackage{graphicx}
\makeatletter
\newsavebox\pandoc@box
\newcommand*\pandocbounded[1]{% scales image to fit in text height/width
  \sbox\pandoc@box{#1}%
  \Gscale@div\@tempa{\textheight}{\dimexpr\ht\pandoc@box+\dp\pandoc@box\relax}%
  \Gscale@div\@tempb{\linewidth}{\wd\pandoc@box}%
  \ifdim\@tempb\p@<\@tempa\p@\let\@tempa\@tempb\fi% select the smaller of both
  \ifdim\@tempa\p@<\p@\scalebox{\@tempa}{\usebox\pandoc@box}%
  \else\usebox{\pandoc@box}%
  \fi%
}
% Set default figure placement to htbp
\def\fps@figure{htbp}
\makeatother

\usepackage{fvextra}
\DefineVerbatimEnvironment{Highlighting}{Verbatim}{breaklines,commandchars=\\\{\}}
\areaset[0.50in]{5.5in}{9in}
\makeatletter
\@ifpackageloaded{tcolorbox}{}{\usepackage[skins,breakable]{tcolorbox}}
\@ifpackageloaded{fontawesome5}{}{\usepackage{fontawesome5}}
\definecolor{quarto-callout-color}{HTML}{909090}
\definecolor{quarto-callout-note-color}{HTML}{0758E5}
\definecolor{quarto-callout-important-color}{HTML}{CC1914}
\definecolor{quarto-callout-warning-color}{HTML}{EB9113}
\definecolor{quarto-callout-tip-color}{HTML}{00A047}
\definecolor{quarto-callout-caution-color}{HTML}{FC5300}
\definecolor{quarto-callout-color-frame}{HTML}{acacac}
\definecolor{quarto-callout-note-color-frame}{HTML}{4582ec}
\definecolor{quarto-callout-important-color-frame}{HTML}{d9534f}
\definecolor{quarto-callout-warning-color-frame}{HTML}{f0ad4e}
\definecolor{quarto-callout-tip-color-frame}{HTML}{02b875}
\definecolor{quarto-callout-caution-color-frame}{HTML}{fd7e14}
\makeatother
\makeatletter
\@ifpackageloaded{bookmark}{}{\usepackage{bookmark}}
\makeatother
\makeatletter
\@ifpackageloaded{caption}{}{\usepackage{caption}}
\AtBeginDocument{%
\ifdefined\contentsname
  \renewcommand*\contentsname{Table of contents}
\else
  \newcommand\contentsname{Table of contents}
\fi
\ifdefined\listfigurename
  \renewcommand*\listfigurename{List of Figures}
\else
  \newcommand\listfigurename{List of Figures}
\fi
\ifdefined\listtablename
  \renewcommand*\listtablename{List of Tables}
\else
  \newcommand\listtablename{List of Tables}
\fi
\ifdefined\figurename
  \renewcommand*\figurename{Figure}
\else
  \newcommand\figurename{Figure}
\fi
\ifdefined\tablename
  \renewcommand*\tablename{Table}
\else
  \newcommand\tablename{Table}
\fi
}
\@ifpackageloaded{float}{}{\usepackage{float}}
\floatstyle{ruled}
\@ifundefined{c@chapter}{\newfloat{codelisting}{h}{lop}}{\newfloat{codelisting}{h}{lop}[chapter]}
\floatname{codelisting}{Listing}
\newcommand*\listoflistings{\listof{codelisting}{List of Listings}}
\makeatother
\makeatletter
\usepackage{pdflscape}
\makeatother
\makeatletter
\makeatother
\makeatletter
\@ifpackageloaded{caption}{}{\usepackage{caption}}
\@ifpackageloaded{subcaption}{}{\usepackage{subcaption}}
\makeatother

\usepackage{bookmark}

\IfFileExists{xurl.sty}{\usepackage{xurl}}{} % add URL line breaks if available
\urlstyle{same} % disable monospaced font for URLs
\hypersetup{
  pdftitle={The Book of OHDSI 2nd Edition},
  hidelinks,
  pdfcreator={LaTeX via pandoc}}


\title{The Book of OHDSI 2nd Edition}
\usepackage{etoolbox}
\makeatletter
\providecommand{\subtitle}[1]{% add subtitle to \maketitle
  \apptocmd{\@title}{\par {\large #1 \par}}{}{}
}
\makeatother
\subtitle{Observational Health Data Sciences and Informatics}
\author{}
\date{2026-04-11}

\begin{document}
\frontmatter
\maketitle

\RecustomVerbatimEnvironment{verbatim}{Verbatim}{
   showspaces = false,
   showtabs = false,
   breaksymbolleft={},
   breaklines
   % Note: setting commandchars=\\\{\} here will cause an error
}

\renewcommand*\contentsname{Table of contents}
{
\setcounter{tocdepth}{2}
\tableofcontents
}

\mainmatter
\bookmarksetup{startatroot}

\chapter*{Welcome}\label{sec-welcome}
\addcontentsline{toc}{chapter}{Welcome}

\markboth{Welcome}{Welcome}

Welcome to 2nd edition of
\href{https://ohdsi.github.io/BookOfOhdsi-2ndEdition/}{\emph{The Book of
OHDSI}}. Here you will find resources for conducting research with the
\href{https://ohdsi.org/}{Observational Health Data Sciences Initiative}
(OHDSI) and the
\href{https://www.ohdsi.org/data-standardization/}{Observational Medical
Outcomes Partnership} CDM (OMOP).

Sincerely,

\emph{The BOO Editorial Committee}

\section*{How to Contribute}\label{sec-welcome-contribute}
\addcontentsline{toc}{section}{How to Contribute}

\markright{How to Contribute}

We welcome suggestions, edits, and larger contributions to this guide.
This book is typeset in Markdown, rendered with
\href{https://quarto.org/}{Quarto}, and hosted on GitHub. For errors or
requests, please submit an
\href{https://docs.github.com/en/issues/tracking-your-work-with-issues/about-issues}{Issue}
to the book's
\href{https://github.com/OHDSI/BookOfOhdsi-2ndEdition/issues}{issue
tracker}. To make larger or direct contributions, please make a
\href{https://github.com/OHDSI/BookOfOhdsi-2ndEdition/pulls}{pull
request} using the standard
\href{https://docs.github.com/en/get-started/quickstart/contributing-to-projects}{GitHub
workflow}.

If you would like to contribute but are unfamiliar with any of these
technologies, please feel free to email QQQ with comments and
suggestions for changes.

\section*{How to Cite this Work}\label{sec-welcome-cite}
\addcontentsline{toc}{section}{How to Cite this Work}

\markright{How to Cite this Work}

O'Neil ST, Beasley W, Loomba J, Patrick S, Wilkins KJ, Crowley KM.,
Anzalone, AJ (Eds.) (2023). \emph{The Researcher's Guide to N3C: A
National Resource for Analyzing Real-World Health Data}. DOI:
\href{https://zenodo.org/record/7749367\#.ZEgPD3bMJmM}{10.5281/zenodo.7749367}

Editorial Committee:

\begin{itemize}
\tightlist
\item
  \href{https://orcid.org/0000-0002-3641-055X}{Christian Reich}: March
  2024 - present
\item
  \href{https://orcid.org/0000-0002-5613-5006}{William H. Beasley}:
  March 2024 - present
\item
  (Add more names \& probably remove me/Will)
\end{itemize}

\section*{Licensing}\label{sec-welcome-licensing}
\addcontentsline{toc}{section}{Licensing}

\markright{Licensing}

This book is licensed under the
\href{https://creativecommons.org/licenses/by-nd/4.0/}{Creative Commons
Attribution-NoDerivatives 4.0}. Individual chapters are as well, unless
otherwise noted.

\section*{Funding}\label{sec-welcome-funding}
\addcontentsline{toc}{section}{Funding}

\markright{Funding}

This content is solely the responsibility of the authors and does not
necessarily represent the official views of QQQ.

\part{OHDSI}

\chapter{The OHDSI Community}\label{sec-ohdsi-community}

TODO- write abstract

\hfill\break

\textbf{Chapter Leads}: George Hripcsak, Patrick Ryan

\begin{tcolorbox}[enhanced jigsaw, opacitybacktitle=0.6, coltitle=black, colback=white, bottomrule=.15mm, bottomtitle=1mm, title=\textcolor{quarto-callout-note-color}{\faInfo}\hspace{0.5em}{Note}, titlerule=0mm, rightrule=.15mm, colbacktitle=quarto-callout-note-color!10!white, colframe=quarto-callout-note-color-frame, toptitle=1mm, arc=.35mm, left=2mm, breakable, opacityback=0, leftrule=.75mm, toprule=.15mm]

This is page is currently a stub. The chapter is being written in the
OHDSI Teams directory. When the draft is complete, it will be converted
to markdown and moved to this file.

Author Resources (requires an OHDSI Teams account):

\begin{itemize}
\tightlist
\item
  \href{https://ohdsiorg.sharepoint.com/sites/Workgroup-EducationWorkingGroup/Shared\%20Documents/Forms/AllItems.aspx?id=\%2Fsites\%2FWorkgroup\%2DEducationWorkingGroup\%2FShared\%20Documents\%2FSection\%20I\%20\%2D\%20Propaganda\%2FChapter\%201\%20The\%20OHDSI\%20Community&viewid=05fec2cc\%2Dec8a\%2D4d04\%2Db565\%2Dcf1289b96f67}{Chapter
  Directory}
\item
  \href{https://ohdsiorg.sharepoint.com/:x:/r/sites/Workgroup-EducationWorkingGroup/_layouts/15/Doc2.aspx?action=edit&sourcedoc=\%7B1fa31e39-1c5f-4918-b878-609ebd9810b3\%7D&wdOrigin=TEAMS-WEB.teamsSdk_ns.rwc&wdExp=TEAMS-TREATMENT&wdhostclicktime=1748104477731&web=1}{Book
  Layout}
\item
  \href{https://ohdsiorg.sharepoint.com/sites/Workgroup-EducationWorkingGroup/Shared\%20Documents/Forms/AllItems.aspx?viewid=05fec2cc\%2Dec8a\%2D4d04\%2Db565\%2Dcf1289b96f67}{Education
  Working Group SharePoint Drive}
\end{itemize}

Public Resources:

\begin{itemize}
\tightlist
\item
  \href{https://ohdsi.github.io/TheBookOfOhdsi/}{Book of OHDSI, Edition
  1}
\item
  \href{https://github.com/OHDSI/TheBookOfOhdsi}{Source Code} for Book
  of OHDSI, Edition 1
\item
  \href{https://ohdsi.org/}{OHDSI Home Page}
\end{itemize}

\end{tcolorbox}

\section{The Journey from Data to
Evidence}\label{sec-ohdsi-community-journey}

\section{OMOP}\label{sec-ohdsi-community-omop}

\section{Adoption}\label{sec-ohdsi-community-adoption}

\section{Stakeholders}\label{sec-ohdsi-community-stakeholders}

\section{Global Diversity}\label{sec-ohdsi-community-global}

\section{Summary}\label{sec-ohdsi-community-summary}

\chapter{OHDSI Principles}\label{sec-ohdsi-principles}

TODO- write abstract

\hfill\break

\textbf{Chapter Leads}: tbc

\begin{tcolorbox}[enhanced jigsaw, opacitybacktitle=0.6, coltitle=black, colback=white, bottomrule=.15mm, bottomtitle=1mm, title=\textcolor{quarto-callout-note-color}{\faInfo}\hspace{0.5em}{Note}, titlerule=0mm, rightrule=.15mm, colbacktitle=quarto-callout-note-color!10!white, colframe=quarto-callout-note-color-frame, toptitle=1mm, arc=.35mm, left=2mm, breakable, opacityback=0, leftrule=.75mm, toprule=.15mm]

This is page is currently a stub. The chapter is being written in the
OHDSI Teams directory. When the draft is complete, it will be converted
to markdown and moved to this file.

Author Resources (requires an OHDSI Teams account):

\begin{itemize}
\tightlist
\item
  \href{https://ohdsiorg.sharepoint.com/sites/Workgroup-EducationWorkingGroup/Shared\%20Documents/Forms/AllItems.aspx?id=\%2Fsites\%2FWorkgroup\%2DEducationWorkingGroup\%2FShared\%20Documents\%2FSection\%20I\%20\%2D\%20Propaganda\%2FChapter\%202\%20OHDSI\%20Principles&viewid=05fec2cc\%2Dec8a\%2D4d04\%2Db565\%2Dcf1289b96f67}{Chapter
  Directory}
\item
  \href{https://ohdsiorg.sharepoint.com/:x:/r/sites/Workgroup-EducationWorkingGroup/_layouts/15/Doc2.aspx?action=edit&sourcedoc=\%7B1fa31e39-1c5f-4918-b878-609ebd9810b3\%7D&wdOrigin=TEAMS-WEB.teamsSdk_ns.rwc&wdExp=TEAMS-TREATMENT&wdhostclicktime=1748104477731&web=1}{Book
  Layout}
\item
  \href{https://ohdsiorg.sharepoint.com/sites/Workgroup-EducationWorkingGroup/Shared\%20Documents/Forms/AllItems.aspx?viewid=05fec2cc\%2Dec8a\%2D4d04\%2Db565\%2Dcf1289b96f67}{Education
  Working Group SharePoint Drive}
\end{itemize}

Public Resources:

\begin{itemize}
\tightlist
\item
  \href{https://ohdsi.github.io/TheBookOfOhdsi/}{Book of OHDSI, Edition
  1}
\item
  \href{https://github.com/OHDSI/TheBookOfOhdsi}{Source Code} for Book
  of OHDSI, Edition 1
\item
  \href{https://ohdsi.org/}{OHDSI Home Page}
\end{itemize}

\end{tcolorbox}

\section{Open Science}\label{sec-ohdsi-principles-science}

\section{Open Standards}\label{sec-ohdsi-principles-standards}

\section{Open Source}\label{sec-ohdsi-principles-source}

\section{Open Data}\label{sec-ohdsi-principles-data}

\section{Open Discourse}\label{sec-ohdsi-principles-discource}

\section{Collaboration}\label{sec-ohdsi-principles-collaboration}

\chapter{Where to Begin}\label{sec-ohdsi-begin}

TODO- write abstract

\hfill\break

\textbf{Chapter Leads}: Kristin Kostka

\begin{tcolorbox}[enhanced jigsaw, opacitybacktitle=0.6, coltitle=black, colback=white, bottomrule=.15mm, bottomtitle=1mm, title=\textcolor{quarto-callout-note-color}{\faInfo}\hspace{0.5em}{Note}, titlerule=0mm, rightrule=.15mm, colbacktitle=quarto-callout-note-color!10!white, colframe=quarto-callout-note-color-frame, toptitle=1mm, arc=.35mm, left=2mm, breakable, opacityback=0, leftrule=.75mm, toprule=.15mm]

This is page is currently a stub. The chapter is being written in the
OHDSI Teams directory. When the draft is complete, it will be converted
to markdown and moved to this file.

Author Resources (requires an OHDSI Teams account):

\begin{itemize}
\tightlist
\item
  \href{https://ohdsiorg.sharepoint.com/sites/Workgroup-EducationWorkingGroup/Shared\%20Documents/Forms/AllItems.aspx?id=\%2Fsites\%2FWorkgroup\%2DEducationWorkingGroup\%2FShared\%20Documents\%2FSection\%20I\%20\%2D\%20Propaganda\%2FChapter\%203\%20Where\%20to\%20Begin&viewid=05fec2cc\%2Dec8a\%2D4d04\%2Db565\%2Dcf1289b96f67}{Chapter
  Directory}
\item
  \href{https://ohdsiorg.sharepoint.com/:x:/r/sites/Workgroup-EducationWorkingGroup/_layouts/15/Doc2.aspx?action=edit&sourcedoc=\%7B1fa31e39-1c5f-4918-b878-609ebd9810b3\%7D&wdOrigin=TEAMS-WEB.teamsSdk_ns.rwc&wdExp=TEAMS-TREATMENT&wdhostclicktime=1748104477731&web=1}{Book
  Layout}
\item
  \href{https://ohdsiorg.sharepoint.com/sites/Workgroup-EducationWorkingGroup/Shared\%20Documents/Forms/AllItems.aspx?viewid=05fec2cc\%2Dec8a\%2D4d04\%2Db565\%2Dcf1289b96f67}{Education
  Working Group SharePoint Drive}
\end{itemize}

Public Resources:

\begin{itemize}
\tightlist
\item
  \href{https://ohdsi.github.io/TheBookOfOhdsi/}{Book of OHDSI, Edition
  1}
\item
  \href{https://github.com/OHDSI/TheBookOfOhdsi}{Source Code} for Book
  of OHDSI, Edition 1
\item
  \href{https://ohdsi.org/}{OHDSI Home Page}
\end{itemize}

\end{tcolorbox}

\section{Join the Journey}\label{sec-ohdsi-begin-join}

\section{Where You Fit In}\label{sec-ohdsi-begin-fit}

\section{Summary}\label{sec-ohdsi-begin-summary}

\part{Uniform Data Representation}

\chapter{The Common Data Model}\label{sec-representation-cdm}

TODO- write abstract

\hfill\break

\textbf{Chapter Leads}: Clair Blacketer

\begin{tcolorbox}[enhanced jigsaw, opacitybacktitle=0.6, coltitle=black, colback=white, bottomrule=.15mm, bottomtitle=1mm, title=\textcolor{quarto-callout-note-color}{\faInfo}\hspace{0.5em}{Note}, titlerule=0mm, rightrule=.15mm, colbacktitle=quarto-callout-note-color!10!white, colframe=quarto-callout-note-color-frame, toptitle=1mm, arc=.35mm, left=2mm, breakable, opacityback=0, leftrule=.75mm, toprule=.15mm]

This is page is currently a stub. The chapter is being written in the
OHDSI Teams directory. When the draft is complete, it will be converted
to markdown and moved to this file.

Author Resources (requires an OHDSI Teams account):

\begin{itemize}
\tightlist
\item
  \href{https://ohdsiorg.sharepoint.com/sites/Workgroup-EducationWorkingGroup/Shared\%20Documents/Forms/AllItems.aspx?id=\%2Fsites\%2FWorkgroup\%2DEducationWorkingGroup\%2FShared\%20Documents\%2FSection\%20II\%20\%2D\%20Data\%2FChapter\%204\%20\%2D\%20Common\%20Data\%20Model&viewid=05fec2cc\%2Dec8a\%2D4d04\%2Db565\%2Dcf1289b96f67}{Chapter
  4 Directory}
\item
  \href{https://ohdsiorg.sharepoint.com/:x:/r/sites/Workgroup-EducationWorkingGroup/_layouts/15/Doc2.aspx?action=edit&sourcedoc=\%7B1fa31e39-1c5f-4918-b878-609ebd9810b3\%7D&wdOrigin=TEAMS-WEB.teamsSdk_ns.rwc&wdExp=TEAMS-TREATMENT&wdhostclicktime=1748104477731&web=1}{Book
  Layout}
\item
  \href{https://ohdsiorg.sharepoint.com/sites/Workgroup-EducationWorkingGroup/Shared\%20Documents/Forms/AllItems.aspx?viewid=05fec2cc\%2Dec8a\%2D4d04\%2Db565\%2Dcf1289b96f67}{Education
  Working Group SharePoint Drive}
\end{itemize}

Public Resources:

\begin{itemize}
\tightlist
\item
  \href{https://ohdsi.github.io/TheBookOfOhdsi/}{Book of OHDSI, Edition
  1}
\item
  \href{https://github.com/OHDSI/TheBookOfOhdsi}{Source Code} for Book
  of OHDSI, Edition 1
\item
  \href{https://ohdsi.org/}{OHDSI Home Page}
\end{itemize}

\end{tcolorbox}

\chapter{Standardized
Vocabularies}\label{sec-representation-vocabularies}

TODO- write abstract

\hfill\break

\textbf{Chapter Leads}: Anna Ostropolets

\begin{tcolorbox}[enhanced jigsaw, opacitybacktitle=0.6, coltitle=black, colback=white, bottomrule=.15mm, bottomtitle=1mm, title=\textcolor{quarto-callout-note-color}{\faInfo}\hspace{0.5em}{Note}, titlerule=0mm, rightrule=.15mm, colbacktitle=quarto-callout-note-color!10!white, colframe=quarto-callout-note-color-frame, toptitle=1mm, arc=.35mm, left=2mm, breakable, opacityback=0, leftrule=.75mm, toprule=.15mm]

This is page is currently a stub. The chapter is being written in the
OHDSI Teams directory. When the draft is complete, it will be converted
to markdown and moved to this file.

Author Resources (requires an OHDSI Teams account):

\begin{itemize}
\tightlist
\item
  \href{https://ohdsiorg.sharepoint.com/sites/Workgroup-EducationWorkingGroup/Shared\%20Documents/Forms/AllItems.aspx?id=\%2Fsites\%2FWorkgroup\%2DEducationWorkingGroup\%2FShared\%20Documents\%2FSection\%20II\%20\%2D\%20Data\%2FChapter\%205\%20\%2D\%20Standardized\%20Vocabularies&viewid=05fec2cc\%2Dec8a\%2D4d04\%2Db565\%2Dcf1289b96f67}{Chapter
  5 Directory}
\item
  \href{https://ohdsiorg.sharepoint.com/:x:/r/sites/Workgroup-EducationWorkingGroup/_layouts/15/Doc2.aspx?action=edit&sourcedoc=\%7B1fa31e39-1c5f-4918-b878-609ebd9810b3\%7D&wdOrigin=TEAMS-WEB.teamsSdk_ns.rwc&wdExp=TEAMS-TREATMENT&wdhostclicktime=1748104477731&web=1}{Book
  Layout}
\item
  \href{https://ohdsiorg.sharepoint.com/sites/Workgroup-EducationWorkingGroup/Shared\%20Documents/Forms/AllItems.aspx?viewid=05fec2cc\%2Dec8a\%2D4d04\%2Db565\%2Dcf1289b96f67}{Education
  Working Group SharePoint Drive}
\end{itemize}

Public Resources:

\begin{itemize}
\tightlist
\item
  \href{https://ohdsi.github.io/TheBookOfOhdsi/}{Book of OHDSI, Edition
  1}
\item
  \href{https://github.com/OHDSI/TheBookOfOhdsi}{Source Code} for Book
  of OHDSI, Edition 1
\item
  \href{https://ohdsi.org/}{OHDSI Home Page}
\end{itemize}

\end{tcolorbox}

The OHDSI Standardized Vocabularies are a foundational part of the OHDSI
research network, and an integral part of the Common Data Model (CDM)
(1). They allow standardization of methods, definitions, and results by
defining the content of the data, paving the way for true remote (behind
the firewall) network research and analytics. Usually, finding and
interpreting the content of observational healthcare data, whether it is
structured data using coding schemes or laid down in free text, is
passed all the way through to the researcher, who is faced with a myriad
of different ways to describe clinical events. OHDSI requires
harmonization not only to a standardized format, but also to a rigorous
standard content.

The Vocabularies are a collection of public standard vocabularies and
terminologies used in the network, which we consolidate from their
different original formats and life-cycle conventions into the CDM table
structure.~The system is dynamic, it evolves with frequent source
vocabulary updates, deprecations, and concept replacements, all of which
are version-controlled and available via ATHENA - OHDSI's vocabulary
distribution platform (2). Vocabularies support phenotyping, covariate
construction, large-scale analytics and result reporting and is a
product of community effort: Vocabulary Team maintains and improves
vocabularies according to the roadmap (3), vocabulary stewards maintain
individual vocabularies (4), various Workgroups coordinate efforts for
special use cases such as device or vaccine harmonization and
contributors across the community add their vocabulary content and
improve existing, particularly through enhancing mappings (5).

In this chapter we first describe the main principles of the
Standardized Vocabularies and processes. We will walk through
Vocabularies components and some typical situations, all of which are
necessary to understand and utilize this foundational resource. We also
point out where the support of the community is required to continuously
improve it and describe special situations that require further
development.

\section{Why Vocabularies, and Why
Standardizing}\label{why-vocabularies-and-why-standardizing}

Medical vocabularies go back to the Bills of Mortality in medieval
London to manage outbreaks of the plague and other diseases
(Figure~\ref{fig-representation-vocabularies-010-london}).

\begin{figure}[H]

\centering{

\pandocbounded{\includegraphics[keepaspectratio]{representation/images/vocabularies/fig-representation-vocabularies-010-london.jpg}}

}

\caption{\label{fig-representation-vocabularies-010-london}1660 London
Bill of Mortality, showing the cause of death for deceased inhabitants
using a classification system of 62 diseases known at the time.}

\end{figure}%

Since then, the classifications have greatly expanded in size and
complexity and spread into other aspects of healthcare, such as
procedures and services, drugs, medical devices, etc. The main
principles have remained the same: they are controlled vocabularies,
terminologies, hierarchies, or ontologies that some healthcare
communities agree upon for the purpose of capturing, classifying, and
analyzing patient data. Many of these vocabularies are maintained by
public and government agencies with a long-term mandate for doing so.
For example, the World Health Organization (WHO) produces the
International Classification of Disease (ICD) with the recent addition
of its 11th revision (ICD11). Local governments create country-specific
versions, such as ICD10CM (USA), ICD10GM (Germany), etc. Governments
also control the marketing and sale of drugs and maintain national
repositories of such certified drugs. Vocabularies are also used in the
private sector, either as commercial products or for internal use, such
as electronic health record (EHR) systems or for medical insurance claim
reporting.

As a result, each country, region, healthcare system and institution
tend to have their own classifications that would most likely only be
relevant where it is used. This myriad of vocabularies prevents
interoperability of the systems they are used in. Standardization is the
key that enables patient data exchange, unlocks health data analysis on
a global level, and allows systematic and standardized research,
including performance characterization and quality assessment. To
address the interoperability problem, multinational organizations have
sprung up and started creating broad standards, such as the Standard
Nomenclature of Medicine (SNOMED) and Logical Observation Identifiers
Names and Codes (LOINC). In the US, the Health IT Standards Committee
(HITAC) recommends the use of SNOMED, LOINC, and the drug vocabulary
RxNorm as standards to the National Coordinator for Health IT (ONC) for
use in a common platform for nationwide health information exchange
across diverse entities.

OHDSI developed the OMOP CDM, a global standard for observational
research. As part of the CDM, the OHDSI Standardized Vocabularies are
available for two main purposes:

\begin{itemize}
\item
  Common repository of all vocabularies used in the community
\item
  Standardization and mapping for use in research
\end{itemize}

The Standardized Vocabularies are available to the community free of
charge and~\textbf{must be used}~for OMOP CDM instance \textbf{as its
mandatory reference table}. It is crucial to use the most recent version
of the Vocabularies and continuously incorporate new versions in the ETL
as Vocabularies changes and shifts impact common research tasks (6).

\subsection{Vocabularies Use Cases and
Users}\label{vocabularies-use-cases-and-users}

OHDSI Vocabularies are different from other ontology systems, such as
the Unified Medical Language System (UMLS) (7) and the difference stems
from the main use case of evidence generation that OHDSI supports. Both
UMLS and OHDSI aggregate relationships from source vocabularies. UMLS
provides crosswalks among vocabularies with various degrees of fidelity
and such crosswalks can be incomplete or ambiguous. OHDSI curates
mappings and selects high-quality ones for the official ``Maps to''
relationships from source vocabularies to a single reference standard,
ensuring that all data sources speak the same language. For example,
``Atrial fibrillation'' from ICD-9, Read, MeSH, etc., are all mapped to
a single SNOMED concept in the Condition domain. UMLS, in contrast,
groups synonymous terms under a CUI but does not designate one as ``the
code to use'' serving as a translation table and not enforcing a single
vocabulary for data encoding. UMLS contains many international
vocabularies, but historically it has had a strong U.S. focus, and some
content can lag in updates. OHDSI Vocabularies explicitly integrate both
US and non-US coding systems and even create new standard concepts for
non-US use cases to achieve global coverage. For example, US drugs are
covered in RxNorm that we import, international drugs are covered in
RxNorm Extension that we create de-novo and both of them are integrated
with Anatomic Therapeutic Classification (ATC) (8,9). OHDSI Vocabularies
are optimized for standardized analytics, offering open-access,
harmonized coverage of both U.S. and international terminologies to
enable consistent, reproducible studies across institutions and
countries.

Vocabularies are centered around generating real-world evidence from
observational studies and are mostly used for two groups of tasks: ETL
of data to OMOP CDM and subsequent research on converted data. If you
are a data engineer/ETL developer, the most relevant information is how
to use correct source-to-standard mappings and populate both standard
and source concept ID fields appropriately. Additionally, you need to
know how to track vocabulary changes adopt ETL accordingly. If you are a
researcher, the most relevant information is how to use vocabularies to
find relevant codes for concept sets and features, use hierarchies, and
examine mappings.

\subsection{Access to the Standardized
Vocabularies}\label{access-to-the-standardized-vocabularies}

The OHDSI Standardized Vocabularies are distributed via ATHENA (2), a
web-based platform for browsing and downloading vocabulary data. You can
use it to search, explore, and filter vocabularies by domain, concept
class, vocabulary source, standard status and validity. You can select
relevant vocabularies and download a pre-packaged vocabulary bundle,
ready for loading into a local OMOP CDM instance.

To download a zip file with all Standardized Vocabularies tables, select
all the vocabularies you need for your OMOP CDM. Vocabularies with
Standard Concepts and very common usage are preselected. Add
vocabularies that are used in your source data. Vocabularies that are
proprietary have no select button. Click on the ``License required''
button to incorporate a licensed-required vocabulary into your list. The
Vocabulary Team will contact you and request you demonstrate your
license or help you connect to the right body to obtain one.

Each vocabulary download includes a ZIP file containing a standard set
of CSV files, which can be loaded into your database using standard SQL
scripts or programmatically. You will also need to re-constitute names
of CPT4 codes as per our use agreement (10).

The VOCABULARY.csv file contains the version and release date metadata
for each vocabulary, which should be recorded to ensure reproducibility
in analyses and network studies. When updating to a newer vocabulary
version, we recommend reviewing the changes in concept definitions,
domain assignments, mappings, and deprecated concepts to ensure that
downstream data and cohort definitions remain valid (6).

You can also select a specific vocabulary release different from the
current release or download a file that contains the delta between two
given releases.

\section{Vocabularies Process and
Governance}\label{vocabularies-process-and-governance}

\subsection{Building the Standardized Vocabularies and Vocabularies
Principles}\label{building-the-standardized-vocabularies-and-vocabularies-principles}

All vocabularies of the Standardized Vocabularies are consolidated into
the same common format: \texttt{CONCEPT},
\texttt{CONCEPT\_RELATIONSHIP}, \texttt{CONCEPT\_ANCESTOR},
\texttt{CONCEPT\_SYNONYM}, and supporting reference files such as
VOCABULARY, DOMAIN, CONCEPT\_CLASS, and RELATIONSHIP. This relieves the
researchers from having to understand and handle multiple different
formats and life-cycle conventions of the originating vocabularies.
\hl{}

OHDSI generally prefers adopting existing vocabularies, rather than
de-novo construction, because (i) many vocabularies have already been
utilized in observational data in the community, and (ii) construction
and maintenance of vocabularies is complex and requires the input of
many stakeholders over long periods of time to mature. For this reason,
dedicated organizations provide vocabularies, which undergo a life cycle
of generation, deprecation, merging, and splitting. Currently, OHDSI
only produces internal administrative vocabularies like Type Concepts
(for example,~condition type concepts) as well as several other
vocabularies to cover areas with existing gaps: RxNorm Extension to
cover drugs that are only used outside the United States, OMOP
Investigational Drugs for investigational drugs, Cancer Modifier for
cancer measurements, and OMOP Extension for miscellaneous gaps. There
are other community-driven efforts, such as GIS Vocabulary Package (11).

All vocabularies go through several common stages upon refresh: staging
or harmonization to a common table structure, normalization and creation
of crosswalks, integration with other vocabularies and release (12). All
steps are accompanied by a set of quality assurance and control
procedures, both automated and human-curated (13).

OHDSI Vocabularies follow twelve principles (14). Among others,
Vocabularies focus on and support OHDSI \textbf{use case} of generating
new evidence. They meant to be~\textbf{comprehensive,} that is~there are
enough concepts to cover any event relevant to the patient's healthcare
experience (e.g.,~conditions, procedures, exposures to drug, etc.) as
well as some of the administrative information of the healthcare system
(e.g.,~visits, care sites, etc.). They strive to have \textbf{unique
standard concept}, where for each Clinical Entity there is only one
concept representing it, called the Standard Concept. Other equivalent
or similar concepts are designated non-Standard and mapped to the
Standard ones. Moreover, such concepts should be stated as fact, no
negations of facts, no reference to the past, and no flavors of null
(unknown, not reported, etc.).

\subsection{Vocabularies Governance, Roadmap and Role of the
Community}\label{vocabularies-governance-roadmap-and-role-of-the-community}

OHDSI Vocabularies work and processes are governed by the OHDSI Central
Coordinating Center's body, Vocabulary Committee, which includes
representatives from across the OHDSI community and helps set priorities
for maintenance, content expansion, and quality improvement. Committee's
work is based on the landscape assessment conducted in 2023 (15). As a
part of this work, it approves the roadmap for bi-annual releases of the
Vocabularies.

Release happens in February and August and is accompanied by detailed
roadmap updates and release notes describing the changes (16). Each
release note describes (1) newly added vocabularies, (2) concepts and/or
mappings newly added to the existing vocabularies, (4) changes in
mappings, domains, status of the concepts as well as detailed
description of the actions performed for specific vocabularies.
Additionally, the release notes contain the artifacts not available
through Athena: pack content, SSSOM-compatible metadata (17) for
concepts and relationships as well as reports for alignment with
vocabulary principles (August 2024 release).

You can use information about releases and the roadmap in two ways.
First, you can assess the content of each release and use open source
tools to assess its impact on ETL and research (6,18). If you are a
ETLer or are responsible for your institution's data/OMOP CDM,
\textbf{you should update OMOP CDM instance with the latest version of
the Vocabularies} to benefit from improved coverage, consistency, and
corrected mappings. Vocabularies \textbf{change} a lot. If you are not
updating, you are falling behind in research.

Second, you can use this information to assess if planned activities
meet your needs. You assist in vocabulary maintenance if your vocabulary
is not on the roadmap as a vocabulary steward (19). With each release
stewards from the community refresh and improve their vocabularies, even
if they are not on the roadmap (current list of stewards can be found
here (4)). You can also add your vocabulary, concepts or improve
existing content (mappings, domains) through community contribution (5).

\subsubsection{Community contributions}\label{community-contributions}

The extensive scope of the OMOP Standard Vocabularies poses a challenge
to maintenance and scalability. The OHDSI Vocabulary Team focuses on
core terminologies out of necessity. However, there are plenty of
opportunities for the OHDSI community to assist in vocabulary
maintenance. Examples above regarding improvements in mappings, labels,
synonyms, etc. are welcomed from the community. Working Groups may feel
particular ownership of a domain or specialty area and wish to help
manage the necessary vocabulary. When this happens, the core terminology
management system can be extended through integration of
community-contributed content. Begun in earnest in 2024, this community
and centrally-management vocabulary integration allows for more scalable
contribution and more rapid conceptual gap-filling. A
community-contribution infrastructure has been developed in phases
depending on the complexity of the contribution. Small changes or
additions can be provided using templates.

You can use templates to add a new standard vocabulary, add non-standard
concepts, add mappings, change mappings or concept domains, or propose
upgrading a non-standard concept to standard. Modification of content
requires community approval through the Vocabulary Workgroup. Template
submissions should be completed and ratified two months prior to the
release date (end of June/end of December for August and February
releases, respectively). Instructions for completed templates are
described on the GitHub Wiki (5).

Larger contributions (for example, entire terminologies or drug
catalogues) require staging and integration using a compatible
environment to that used for managing the core terminologies. Examples
of community-staged contributions include the Heme-Onc vocabulary, the
Veterinary vocabulary and the CIEL terminology. For complex
contributions, it is best to have a working group sponsor your request.
You can use the instructions provided on Wiki under Community
Contributions Part II (5). We recommend you talk to the members of the
Vocabulary Workgroup or Team to discuss your specific use case.

\section{OHDSI Vocabularies Structure: Concepts and
Relationships}\label{ohdsi-vocabularies-structure-concepts-and-relationships}

All clinical events in the OMOP CDM are represented as concepts, which
capture the semantic notion of each event. They are the fundamental
building blocks of the data records, making almost all tables fully
normalized with few exceptions. Concepts are stored in the
\texttt{CONCEPT} table
(Figure~\ref{fig-representation-vocabularies-020-standard-mapping}).

\begin{figure}[H]

\centering{

\pandocbounded{\includegraphics[keepaspectratio]{representation/images/vocabularies/fig-representation-vocabularies-020-standard-mapping.png}}

}

\caption{\label{fig-representation-vocabularies-020-standard-mapping}Standard
representation of vocabulary concepts in the OMOP CDM. The example
provided is the \texttt{CONCEPT} table record for the SNOMED code for
Atrial Fibrillation.}

\end{figure}%

\subsection{Concept IDs}\label{concept-ids}

Each concept is assigned a concept ID to be used as a primary key. This
meaningless integer ID, rather than the original code from the source
vocabulary, is used to record data in the CDM event tables via the
foreign key fields. No two concepts (even from different vocabularies)
share the same ID. Conversely, the same source code might appear in
multiple vocabularies, but each distinct concept gets its own ID.

\subsection{Concept Names}\label{concept-names}

Each concept has one name. Names are always in English. They are
imported from the source of the vocabulary. If the source vocabulary has
more than one name, the most expressive (fully specified) is selected
and the remaining ones are stored in the \texttt{CONCEPT\_SYNONYM} table
under the same \texttt{CONCEPT\_ID} key. Non-English names are recorded
in \texttt{CONCEPT\_SYNONYM} as well, with the appropriate language
concept ID in the \texttt{LANGUAGE\_CONCEPT\_ID} field. The name can
only be 255 characters long, which means that very long names get
truncated, and the full-length version recorded as another synonym,
which can hold up to 1000 characters. Tools like Athena and ATLAS use
the concept names and synonyms to let users search for concepts. When
doing analysis, it is often convenient to have the concept names for
interpretability, but analysis logic should use the
\texttt{CONCEPT\_ID}.

\subsection{Domains}\label{domains}

Each concept is assigned a domain in the \texttt{DOMAIN\_ID} field,
which, in contrast to the numerical \texttt{CONCEPT\_ID}, is a short,
case-sensitive, unique alphanumeric ID for the domain. Domains are
OMOP-specific and correspond to the OMOP CDM tables (20). Examples of
such identifiers are ``Condition,'' ``Drug,'' ``Procedure,'' ``Visit,''
``Device,'' ``Specimen,'' etc. Domains also direct to which CDM table
and field a clinical event or event attribute is recorded. For example,
``Atrial fibrillation'' is a clinical finding that would be recorded in
the Condition Occurrence table, so its domain is ``Condition''; a
concept for a lab test (for example, ``Blood glucose measurement'')
would have domain ``Measurement'' and belong in the Measurement table.
Domains are assigned to codes, and a vocabulary can have different
domains: for example, HCPCS, while considered procedure vocabulary, also
has codes with Drug and Observation domains.

The domain heuristic follows the definitions of the domains. These
definitions are derived from the table and field definitions in the CDM
\textbf{\{Chapter 4\}}. The heuristic is not perfect; there are grey
zones (\textbf{\{Section 5.4\}~}''Special Situations''), source
vocabulary shifts, and occasional misassignments. Although domains of
concepts may change, 95\% of the concepts never changed their domain
since Vocabularies' inception (for more information, see Assets in
v20240830 release notes) (16).

\subsection{Vocabularies}\label{vocabularies}

Each vocabulary has a short case-sensitive unique alphanumeric ID, which
generally follows the abbreviated name of the vocabulary, omitting
dashes. For example, ICD-9-CM has the vocabulary ID ``ICD9CM''. As of
2025, over 140 vocabularies are available through ATHENA and follow
different cadence of updates. The source and the version of the
vocabularies is defined in the VOCABULARY reference file and
documentation for individual vocabularies can be found on GitHub (4,16).

\subsection{Concept Classes}\label{concept-classes}

Some vocabularies classify their codes or concepts, denoted through
their case-sensitive unique alphanumerical IDs. For example, SNOMED has
33 such concept classes, which SNOMED refers to as ``semantic tags'':
clinical finding, social context, body structure, etc. These are
vertical divisions of the concepts. Others, such as MedDRA or RxNorm,
have concept classes classifying horizontal levels in their stratified
hierarchies. Vocabularies without any concept classes, such as HCPCS,
use the vocabulary ID as the Concept Class ID.

\begin{longtable}[]{@{}
  >{\raggedright\arraybackslash}p{(\linewidth - 2\tabcolsep) * \real{0.2606}}
  >{\raggedright\arraybackslash}p{(\linewidth - 2\tabcolsep) * \real{0.7394}}@{}}
\caption{Vocabularies with or without horizontal and vertical
sub-classification principles in concept
class.}\label{tbl-part-chapter-shorttableterm}\tabularnewline
\toprule\noalign{}
\begin{minipage}[b]{\linewidth}\raggedright
Concept class subdivision principle
\end{minipage} & \begin{minipage}[b]{\linewidth}\raggedright
Vocabulary
\end{minipage} \\
\midrule\noalign{}
\endfirsthead
\toprule\noalign{}
\begin{minipage}[b]{\linewidth}\raggedright
Concept class subdivision principle
\end{minipage} & \begin{minipage}[b]{\linewidth}\raggedright
Vocabulary
\end{minipage} \\
\midrule\noalign{}
\endhead
\bottomrule\noalign{}
\endlastfoot
Horizontal & All drug vocabularies, CDТ, Episode, HCPCS, HemOnc, ICDs,
MedDRA, OSM, Census \\
Vertical & CIEL, HES Specialty, ICDO3, MeSH, NAACCR, NDFRT, OPCS4,
PCORNET, Plan, PPI, Provider, SNOMED, SPL, UCUM \\
Mixed & CPT4, ISBT, LOINC \\
None & OXMIS, Race, Revenue Code, Sponsor, Supplier, UB04s, Visit \\
\end{longtable}

Horizontal concept classes allow you to determine a specific
hierarchical level. For example, in the drug vocabulary RxNorm, the
concept class ``Ingredient'' defines the top level of the hierarchy. In
the vertical model, members of a concept class can be of any
hierarchical level from the top to the very bottom. Concept class is
mostly a descriptive attribute and helps to filter concepts. For
example, if you only want to select drugs with a specific Brand Name you
can filter to ``Branded Drug'' class.

\subsection{Standard Concepts}\label{standard-concepts}

A Standard Concept is the community-endorsed, canonical representation
of a clinical meaning within the OHDSI Vocabularies. It serves as the
unified semantic identifier for a specific entity (for example,
condition, drug, procedure), regardless of how that entity is expressed
in source vocabularies. Only Standard Concepts are used to populate the
\texttt{CONCEPT\_ID} fields in the CDM, ensuring consistency across
diverse datasets. Standard concepts serve as the target for mappings.
For each clinical entity, one concept from one vocabulary is chosen to
be standard. This becomes the ``hub'' to which all equivalent source
codes are mapped. For example, MESH code D001281, CIEL code 148203,
SNOMED code 49436004, ICD9CM code 427.31 and Read code G573000 all
define ``Atrial fibrillation'' in the condition domain, but only the
SNOMED concept is Standard and represents the condition in the data. The
others are designated non-standard or source concepts and mapped to the
Standard ones. Standard Concepts are indicated through an ``S'' in the
\texttt{STANDARD\_CONCEPT} field. And only these Standard Concepts are
used to record data in the CDM fields ending in \texttt{\_CONCEPT\_ID}.

We rely on well-known reference terminologies for standard terms: SNOMED
CT for conditions, RxNorm and RxNorm Extension for drugs, LOINC and
SNOMED for measurements, etc. Not all concepts in those vocabularies are
necessarily standard. Occasionally, a concept in a standard vocabulary
might be deemed out of scope or duplicative and not used. Conversely,
some concepts from typically non-standard vocabularies might be made
standard if no better alternative exists.

While we strive to align with the unique standard concept principle to
have one Standard concept per semantic entity, duplicates exist. For
example, no deduplication of standard concepts has been performed for
the Device domain. While concept mappings avoid direct collisions, this
dual-standard condition can introduce ambiguity in cohort definition,
concept set construction, and analytic interpretation. This phenomenon
is not a flaw but rather a reflection of ontology convergence in
progress, where two high-quality terminologies independently arrive at
comparable representations of the same clinical reality. OHDSI addresses
these cases through community review, classification logic, and
long-term efforts toward consolidation via concept deprecation,
reclassification, or updated mappings. Until then, such duplications
must be handled with care in concept set design and ETL strategies.

\subsection{Non-Standard Concepts}\label{non-standard-concepts}

Non-standard concepts are not used in standardized analytics, but they
are still part of the Standardized Vocabularies and are often found in
the source data. For that reason, they are also called ``source
concepts''. The conversion of source concepts to Standard Concepts is a
process called ``mapping''.

Some of the non-standard concepts cannot be mapped and are not suitable
for analytic use. Examples of such include terms like ``Not reported'',
``Not specified'', ``Passport number'' and more.

\subsection{Classification Concepts}\label{classification-concepts}

These concepts are not Standard and hence cannot be used to represent
the data. But they are participating in the hierarchy with the Standard
Concepts and can therefore be used to perform hierarchical queries. For
example, querying for all descendants of ATC code prednisolone;systemic
will retrieve the Standard RxNorm concept for prednisolone 5 MG Oral
Tablet
(Figure~\ref{fig-representation-vocabularies-030-standard-hierarchy}).

\begin{landscape}

\begin{figure}[H]

\centering{

\pandocbounded{\includegraphics[keepaspectratio]{representation/images/vocabularies/fig-representation-vocabularies-030-standard-hierarchy.png}}

}

\caption{\label{fig-representation-vocabularies-030-standard-hierarchy}Standard,
non-standard source and classification concepts and their hierarchical
relationships in the drug domain.}

\end{figure}%

\end{landscape}

Classification concepts are marked with a ``C'' in the
\texttt{STANDARD\_CONCEPT} field. Most classification concepts form a
hierarchy along with the standard concepts, and these relationships are
stored in the \texttt{CONCEPT\_ANCESTOR} table.

Classification concepts are vital in enabling concept set expansion via
ancestry traversal. While they cannot be used to populate clinical event
tables directly, they serve as entry points into clinically meaningful
groupings (for example, drug classes or disorder categories). They are
especially powerful when used in cohort definitions or phenotype
algorithms that require aggregation of clinically related Standard
Concepts. For example, selecting the ATC classification concept C09AA
ACE inhibitors will retrieve all Standard RxNorm ingredients and
products mapped as descendants.

Quality and coverage of classification hierarchies vary across domains.
Currently, Drug and Condition domains have mature classification
structures (for example, ATC, MedDRA), while Procedure, Measurement, and
Device domains lack formal classification vocabularies. Caution should
be exercised when interpreting classification-derived hierarchies, as
they may not always reflect clinical practice or data granularity.

5.2.8.1 Standard/non-standard/classification Concept Designation

The choice of concept designation as Standard, non-standard, and
classification is typically done for each domain separately at the
vocabulary level. This is based on the quality of the concepts, the
built-in hierarchy, and the declared purpose of the vocabulary. Also,
not all concepts in a vocabulary are used as Standard Concepts. The
designation is separate for each domain, each concept must be active,
and there might be an order of precedence if more than one concept from
different vocabularies compete for the same meaning. See Table~5.2~for
examples.

\begin{landscape}

\begin{longtable}[]{@{}
  >{\raggedright\arraybackslash}p{(\linewidth - 6\tabcolsep) * \real{0.0915}}
  >{\raggedright\arraybackslash}p{(\linewidth - 6\tabcolsep) * \real{0.3380}}
  >{\raggedright\arraybackslash}p{(\linewidth - 6\tabcolsep) * \real{0.3662}}
  >{\raggedright\arraybackslash}p{(\linewidth - 6\tabcolsep) * \real{0.2042}}@{}}
\caption{List of vocabularies to utilize for
Standard/non-standard/classification concept
assignments.}\label{tbl-part-chapter-shorttableterm}\tabularnewline
\toprule\noalign{}
\begin{minipage}[b]{\linewidth}\raggedright
Domain
\end{minipage} & \begin{minipage}[b]{\linewidth}\raggedright
for Standard Concepts
\end{minipage} & \begin{minipage}[b]{\linewidth}\raggedright
for source concepts
\end{minipage} & \begin{minipage}[b]{\linewidth}\raggedright
for classification concepts
\end{minipage} \\
\midrule\noalign{}
\endfirsthead
\toprule\noalign{}
\begin{minipage}[b]{\linewidth}\raggedright
Domain
\end{minipage} & \begin{minipage}[b]{\linewidth}\raggedright
for Standard Concepts
\end{minipage} & \begin{minipage}[b]{\linewidth}\raggedright
for source concepts
\end{minipage} & \begin{minipage}[b]{\linewidth}\raggedright
for classification concepts
\end{minipage} \\
\midrule\noalign{}
\endhead
\bottomrule\noalign{}
\endlastfoot
Condition & SNOMED, ICDO3 & SNOMED Veterinary & MedDRA \\
Procedure & SNOMED, CPT4, HCPCS, ICD10PCS, ICD9Proc, OPCS4 & SNOMED
Veterinary, HemOnc, NAACCR & None at this point \\
Measurement & SNOMED, LOINC & SNOMED Veterinary, NAACCR, CPT4, HCPCS,
OPCS4, PPI & None at this point \\
Drug & RxNorm, RxNorm Extension, CVX & HCPCS, CPT4, HemOnc, NAACCR &
ATC \\
Device & SNOMED & Others, currently not normalized & None at this
point \\
Observation & SNOMED & Others & None at this point \\
Visit & CMS Place of Service, ABMT, NUCC & SNOMED, HCPCS, CPT4, UB04 &
None at this point \\
\end{longtable}

\end{landscape}

\subsection{Concept Codes}\label{concept-codes}

Concept codes are the identifiers used in the source vocabularies. For
example, ICD9CM or NDC codes are stored in this field, while the OMOP
tables use the concept ID as a foreign key into the \texttt{CONCEPT}
table. The reason is that the name space overlaps across vocabularies,
that is~the same code can exist in different vocabularies with
completely different meanings (Table~5.3).

\begin{landscape}

\begin{longtable}[]{@{}
  >{\raggedleft\arraybackslash}p{(\linewidth - 10\tabcolsep) * \real{0.0945}}
  >{\raggedleft\arraybackslash}p{(\linewidth - 10\tabcolsep) * \real{0.1102}}
  >{\raggedright\arraybackslash}p{(\linewidth - 10\tabcolsep) * \real{0.4331}}
  >{\raggedright\arraybackslash}p{(\linewidth - 10\tabcolsep) * \real{0.1102}}
  >{\raggedright\arraybackslash}p{(\linewidth - 10\tabcolsep) * \real{0.1181}}
  >{\raggedright\arraybackslash}p{(\linewidth - 10\tabcolsep) * \real{0.1339}}@{}}
\caption{Concepts with identical concept code 1001, but different
vocabularies, domains and concept
classes.}\label{tbl-part-chapter-shorttableterm}\tabularnewline
\toprule\noalign{}
\begin{minipage}[b]{\linewidth}\raggedleft
Concept ID
\end{minipage} & \begin{minipage}[b]{\linewidth}\raggedleft
Concept Code
\end{minipage} & \begin{minipage}[b]{\linewidth}\raggedright
ConceptName
\end{minipage} & \begin{minipage}[b]{\linewidth}\raggedright
DomainID
\end{minipage} & \begin{minipage}[b]{\linewidth}\raggedright
VocabularyID
\end{minipage} & \begin{minipage}[b]{\linewidth}\raggedright
ConceptClass
\end{minipage} \\
\midrule\noalign{}
\endfirsthead
\toprule\noalign{}
\begin{minipage}[b]{\linewidth}\raggedleft
Concept ID
\end{minipage} & \begin{minipage}[b]{\linewidth}\raggedleft
Concept Code
\end{minipage} & \begin{minipage}[b]{\linewidth}\raggedright
ConceptName
\end{minipage} & \begin{minipage}[b]{\linewidth}\raggedright
DomainID
\end{minipage} & \begin{minipage}[b]{\linewidth}\raggedright
VocabularyID
\end{minipage} & \begin{minipage}[b]{\linewidth}\raggedright
ConceptClass
\end{minipage} \\
\midrule\noalign{}
\endhead
\bottomrule\noalign{}
\endlastfoot
35803438 & 1001 & Granulocyte colony-stimulating factors & Drug & HemOnc
& Component Class \\
35942070 & 1001 & AJCC TNM Clin T & Measurement & NAACCR & NAACCR
Variable \\
1036059 & 1001 & Antipyrine & Drug & RxNorm & Ingredient \\
38003544 & 1001 & Residential Treatment - Psychiatric & Revenue Code &
Revenue Code & Revenue Code \\
43228317 & 1001 & Aceprometazine maleate & Drug & BDPM & Ingredient \\
45417187 & 1001 & Brompheniramine Maleate, 10 mg/mL injectable solution
& Drug & Multum & Multum \\
45912144 & 1001 & Serum & Specimen & CIEL & Specimen \\
\end{longtable}

\end{landscape}

\texttt{CONCEPT\_CODE} is unique only within a given vocabulary. You
should not join datasets via \texttt{CONCEPT\_CODE} unless constrained
by \texttt{VOCABULARY\_ID}.

In addition, certain vocabularies, such as HCPCS, NDC, and DRG are known
to reuse codes over time, assigning new meanings to previously used
codes. In such cases, Vocabularies differentiate concepts based on
validity dates (\texttt{VALID\_START\_DATE}, \texttt{VALID\_END\_DATE})
and keep the most recent meaning.

Some OMOP-specific vocabularies (for example, Type Concept, Visit)
contain system-generated concept codes rather than real-world codes.
Finally, certain source vocabularies (such as ATC or hierarchical
clinical classifications) embed structural hierarchy into their codes
(ATC G03E vs.~G03EK), meaning that not all \texttt{CONCEPT\_CODE}
matches imply equivalence at the clinical level.

\subsection{Lifecycle}\label{lifecycle}

Vocabularies are rarely permanent corpora with a fixed set of codes.
Instead, codes and concepts are added and get deprecated. The OMOP CDM
is a model to support longitudinal patient data, which means it needs to
support concepts that were used in the past and might no longer be
active, as well as supporting new concepts and placing them into
context. There are three fields in the \texttt{CONCEPT} table that
describe the possible life-cycle statuses: \texttt{VALID\_START\_DATE},
\texttt{VALID\_END\_DATE}, and \texttt{INVALID\_REASON}. Their values
differ depending on the concept life-cycle status:

\begin{itemize}
\item
  \textbf{Active or new concept}

  \begin{itemize}
  \item
    Description: Concept in use.
  \item
    \texttt{VALID\_START\_DATE}: Day of instantiation of concept; if
    that is not known, day of incorporation of concept in Vocabularies;
    if that is not known, 1970-1-1.
  \item
    \texttt{VALID\_END\_DATE}: Set to 2099-12-31 as a convention to
    indicate ``Might become invalid in an undefined future, but active
    right now''.
  \item
    \texttt{INVALID\_REASON}: NULL
  \end{itemize}
\item
  \textbf{Deprecated Concept with no successor}

  \begin{itemize}
  \item
    Description: Concept inactive and cannot be used as Standard.
  \item
    \texttt{VALID\_START\_DATE}: Day of instantiation of concept; if
    that is not known, day of incorporation of concept in Vocabularies;
    if that is not known, 1970-1-1.
  \item
    \texttt{VALID\_END\_DATE}: Day in the past indicating deprecation,
    or if that is not known, day of vocabulary refresh where concept in
    vocabulary went missing or set to inactive.
  \item
    \texttt{INVALID\_REASON}: ``D''
  \end{itemize}
\item
  \textbf{Upgraded Concept with successor}

  \begin{itemize}
  \item
    Description: Concept inactive but has defined successor. These are
    typically concepts which went through de-duplication.
  \item
    \texttt{VALID\_START\_DATE}: Day of instantiation of concept; if
    that is not known, day of incorporation of concept in Vocabularies;
    if that is not known, 1970-1-1.
  \item
    \texttt{VALID\_END\_DATE}: Day in the past indicating an upgrade, or
    if that is not known day of vocabulary refresh where the upgrade was
    included.
  \item
    \texttt{INVALID\_REASON}: ``U''
  \end{itemize}
\item
  \textbf{Reused code for another new concept}

  \begin{itemize}
  \item
    Description: The vocabulary reused the concept code of this
    deprecated concept for a new concept.
  \item
    \texttt{VALID\_START\_DATE}: Day of instantiation of concept; if
    that is not known, day of incorporation of concept in Vocabularies;
    if that is not known, 1970-1-1.
  \item
    \texttt{VALID\_END\_DATE}: Day in the past indicating deprecation,
    or if that is not known day of vocabulary refresh where concept in
    vocabulary went missing or set to inactive.
  \end{itemize}
\end{itemize}

In addition to concept lifecycle management, the
\texttt{CONCEPT\_RELATIONSHIP} table also has lifecycle fields
(\texttt{VALID\_START\_DATE}, \texttt{VALID\_END\_DATE},
\texttt{INVALID\_REASON}) for relationships. Relationships may change
over time independently of the concepts themselves. While all
relationships are versioned in the internal vocabulary system, only
active mappings are included in Athena downloads. Every OMOP CDM
instance should record the vocabulary version (stored in the
\texttt{VOCABULARY} table) used at ETL time to ensure transparency and
reproducibility. Lifecycle management principles apply equally to custom
extensions and community-contributed vocabularies: all new concepts and
mappings must carry valid \texttt{VALID\_START\_DATE} entries and, when
deprecated, clearly marked \texttt{VALID\_END\_DATE} and
\texttt{INVALID\_REASON} values.

\subsection{Relationships}\label{relationships}

Any two concepts can have a defined relationship, regardless of whether
the two concepts belong to the same domain or vocabulary. The nature of
the relationships is indicated in its short case-sensitive unique
alphanumeric ID in the \texttt{RELATIONSHIP\_ID} field of the
\texttt{CONCEPT\_RELATIONSHIP} table. Relationships are symmetrical,
that is~for each relationship an equivalent relationship exists, where
the content of the fields \texttt{CONCEPT\_ID\_1} and
\texttt{CONCEPT\_ID\_2} are swapped, and the \texttt{RELATIONHSIP\_ID}
is changed to its opposite. For example, the ``Maps to'' relationship
has an opposite relationship ``Mapped from.'' Different types of
relationships serve different analytic purposes. ``Maps to'' and
``Mapped from'' support source-to-standard mappings. ``Is a'' and
``Subsumes'' define hierarchical subclass relationships. ``Has
ingredient'' and ``Ingredient of'' structure drug compositions.
``Concept replaced by'' and ``Concept replaces'' handle lifecycle
transitions across deprecated content.

As stated in the previous section, \texttt{CONCEPT\_RELATIONSHIP} table
records also have life-cycle fields \texttt{VALID\_START\_DATE},
\texttt{VALID\_END\_DATE} and \texttt{INVALID\_REASON}. However, only
active records with \texttt{INVALID\_REASON\ =\ NULL} are available
through ATHENA. Inactive relationships are kept for internal processing
only.

The \texttt{RELATIONSHIP} table serves as the reference with the full
list of relationship IDs and their reverse counterparts. It also
specifies two important flags: \texttt{DEFINES\_ANCESTRY}, indicating
whether a relationship should contribute to the
\texttt{CONCEPT\_ANCESTOR} table, and \texttt{IS\_HIERARCHICAL},
indicating whether the relationship encodes a subsumption hierarchy. Not
all relationships define ancestry; only those intended to build domain
hierarchies (for example, ``Is a'') are used to populate
\texttt{CONCEPT\_ANCESTOR}. It is essential to distinguish between
direct relationships (stored in \texttt{CONCEPT\_RELATIONSHIP}) and
inferred multi-level hierarchies (precomputed and stored in
\texttt{CONCEPT\_ANCESTOR}), especially when writing concept set
queries, building phenotypes, or exploring ontology structures.

\subsection{Mapping Relationships}\label{mapping-relationships}

These relationships provide translations from non-standard to Standard
concepts, supported by two relationship ID pairs (Table~5.4).

\begin{longtable}[]{@{}
  >{\raggedright\arraybackslash}p{(\linewidth - 2\tabcolsep) * \real{0.3000}}
  >{\raggedright\arraybackslash}p{(\linewidth - 2\tabcolsep) * \real{0.7000}}@{}}
\caption{Type of mapping
relationships.}\label{tbl-part-chapter-shorttableterm}\tabularnewline
\toprule\noalign{}
\begin{minipage}[b]{\linewidth}\raggedright
Relationship\\
ID pair\strut
\end{minipage} & \begin{minipage}[b]{\linewidth}\raggedright
Purpose
\end{minipage} \\
\midrule\noalign{}
\endfirsthead
\toprule\noalign{}
\begin{minipage}[b]{\linewidth}\raggedright
Relationship\\
ID pair\strut
\end{minipage} & \begin{minipage}[b]{\linewidth}\raggedright
Purpose
\end{minipage} \\
\midrule\noalign{}
\endhead
\bottomrule\noalign{}
\endlastfoot
\begin{minipage}[t]{\linewidth}\raggedright
``Maps to''\\
and\\
``Mapped from''\strut
\end{minipage} & Mapping to Standard Concepts. Standard Concepts are
mapped to themselves, non-standard concepts to Standard Concepts. Most
non-standard and all Standard Concepts have this relationship to a
Standard Concept. The former are stored in *\_SOURCE\_CONCEPT\_ID, and
the latter in the *\_CONCEPT\_ID fields. Classification concepts are not
mapped. \\
\begin{minipage}[t]{\linewidth}\raggedright
``Maps to value''\\
and\\
``Value mapped from''\strut
\end{minipage} & Mapping to a concept that represents a Value to be
placed into the VALUE\_AS\_CONCEPT\_ID fields of the MEASUREMENT and
OBSERVATION tables. \\
\end{longtable}

The purpose of these mapping relationships is to allow a crosswalk
between equivalent concepts to harmonize how clinical events are
represented in the OMOP CDM. This is a main achievement of the
Standardized Vocabularies.

``Equivalent concepts'' means it carries the same meaning, and,
importantly, the hierarchical descendants cover the same semantic space.
If an equivalent concept is not available and the concept is not
Standard, it is still mapped, but to a slightly broader concept
(so-called ``up-hill mappings'' or semantic subsumption\textbf{)}. For
example, ICD10CM W61.51 ``Bitten by goose'' has no equivalent in the
SNOMED vocabulary, which is generally used for standard condition
concepts. Instead, it is mapped to SNOMED 217716004 ``Peck by bird,''
losing the context of the bird being a goose. Up-hill mappings are only
used if the loss of information is considered irrelevant to standard
research use cases.

Some mappings connect a source concept to more than one Standard
Concept. For example, ICD9CM 070.43 ``Hepatitis E with hepatic coma'' is
mapped to both SNOMED 235867002 ``Acute hepatitis E'' as well as SNOMED
72836002 ``Hepatic Coma.'' The reason for this is that the original
source concept is a pre-coordinated combination of two conditions,
hepatitis and coma, meaning that a single code simultaneously encodes
multiple clinical ideas rather than expressing them separately SNOMED
does not have that combination, which results in two records written to
the \texttt{CONDITION\_OCCURRENCE} table for the single ICD9CM record,
one with each mapped Standard Concept.

Relationships ``Maps to value'' have the purpose of splitting of a value
for OMOP CDM tables following an entity-attribute-value (EAV) model
(21). This is typically the case in the following situations:

\begin{itemize}
\item
  Measurements consisting of a test and a result value
\item
  Personal or family disease history
\item
  Allergy to substance
\item
  Need for immunization
\end{itemize}

In these situations, the source concept is a combination of the
attribute (test or history) and the value (test result or disease). The
``Maps to'' relationship maps this source to the attribute concept, and
the ``Maps to value'' to the value concept. See
Figure~\ref{fig-representation-vocabularies-040-one-to-many}~for an
example.

\begin{figure}[H]

\centering{

\pandocbounded{\includegraphics[keepaspectratio]{representation/images/vocabularies/fig-representation-vocabularies-040-one-to-many.png}}

}

\caption{\label{fig-representation-vocabularies-040-one-to-many}One-to-many
mapping between source concept and Standard Concepts. A pre-coordinated
concept is split into two concepts, one of which is the attribute (here
history of clinical finding) and the other one is the value (peptic
ulcer). While ``Maps to'' relationship will map to concepts of the
measurement or observation domains, the ``Maps to value'' concepts have
no domain restriction.}

\end{figure}%

This process represents a form of controlled \textbf{post-coordination}
within OMOP vocabularies: instead of encoding every possible combination
as a new standard concept, the meaning is decomposed into two (or more)
standardized elements that together fully represent the clinical event.
Together, they enable more flexible, semantically rich, and extensible
data modeling. By post-coordinating attribute and value concepts, OHDSI
Standardized Vocabularies avoid uncontrolled growth in the number of
concepts while still allowing detailed, clinically meaningful data
representation and analysis. Analysts must retrieve both the
\texttt{CONCEPT\_ID} and \texttt{VALUE\_AS\_CONCEPT\_ID} fields together
during query building to reconstruct the complete meaning.

Mapping relationships themselves are subject to lifecycle management.
Deprecated mappings (mappings with an \texttt{INVALID\_REASON} other
than NULL) are removed from active ATHENA releases but can impact
longitudinal data or historical cohort definitions if not updated.
Careful management of mapping versioning is crucial during vocabulary
refresh cycles.

When interpreting mappings, users must be aware that not all
source-to-standard mappings imply perfect semantic equivalence. Slight
loss of detail, context shift, or broader aggregation may occur,
particularly in uphill mappings or when representing pre-coordinated
concepts. Analysts and ETL designers should validate mappings in
critical analytic contexts.

\subsection{Hierarchical Relationships and
Hierarchy}\label{hierarchical-relationships-and-hierarchy}

Relationships which indicate a hierarchy are defined through the ``Is
a'' - ``Subsumes'' relationship pair. Hierarchical relationships are
defined such that the child concept has all the attributes of the parent
concept, plus one or more additional attributes or a more precisely
defined attribute. For example, SNOMED 49436004 ``Atrial fibrillation''
is related to SNOMED 17366009 ``Atrial arrhythmia'' through a ``Is a''
relationship. Both concepts have an identical set of attributes except
the type of arrhythmia, which is defined as fibrillation in one but not
the other. Concepts can have more than one parent and more than one
child concept. In this example, SNOMED 49436004 ``Atrial fibrillation''
is also an ``Is a'' to SNOMED 40593004 ``Fibrillation.''

Within a domain, and in some cases across domains, standard and
classification concepts are organized in a hierarchical structure and
stored in the \texttt{CONCEPT\_ANCESTOR} table. This allows querying and
retrieving concepts and all their hierarchical descendants. These
descendants have the same attributes as their ancestor, but also
additional or more defined ones.

The \texttt{CONCEPT\_ANCESTOR} table is built automatically from the
\texttt{CONCEPT\_RELATIONSHIP} table, traversing all possible concepts
connected through hierarchical relationships. These are the ``Is a'' -
``Subsumes'' pairs
(Figure~\ref{fig-representation-vocabularies-050-ancestor}), and other
relationships connecting hierarchies across vocabularies (``SNOMED -
CPT4 equivalent'', ``RxNorm ingredient of''). The choice whether a
relationship participates in the hierarchy constructor is defined for
each relationship ID by the flag \texttt{DEFINES\_ANCESTRY} in the
\texttt{RELATIONSHIP} reference table. It is important to note that not
all relationships with hierarchical meaning (\texttt{IS\_HIERARCHICAL} =
1) are used for ancestry building; only those with
\texttt{DEFINES\_ANCESTRY} = 1 contribute to \texttt{CONCEPT\_ANCESTOR}.
Relationships such as ``Has FDA approved indication'' or ``Consists of''
are conceptually hierarchical but are excluded from ancestry paths to
preserve clinical rigor.

\begin{figure}[H]

\centering{

\pandocbounded{\includegraphics[keepaspectratio]{representation/images/vocabularies/fig-representation-vocabularies-050-ancestor.png}}

}

\caption{\label{fig-representation-vocabularies-050-ancestor}Hierarchy
of the condition ``Atrial fibrillation''. First degree ancestry is
defined through ``Is a'' and ``Subsumes'' relationships, while all
higher degree relations are inferred and stored in the
\texttt{CONCEPT\_ANCESTOR} table. Each concept is also its own
descendant with both levels of separation equal to 0.}

\end{figure}%

The ancestral degree, or the number of steps between ancestor and
descendant, is captured in the \texttt{MIN\_LEVELS\_OF\_SEPARATION} and
\texttt{MAX\_LEVELS\_OF\_SEPARATION} fields, defining the shortest or
longest possible connection. Not all hierarchical relationships
contribute equally to the levels-of-separation calculation. A step
counted for the degree is determined by the \texttt{IS\_HIERARCHICAL}
flag in the \texttt{RELATIONSHIP} reference table for each relationship
ID.

As of 2025, a high-quality comprehensive hierarchy exists only for two
domains: Drug and Condition. Procedure, Measurement, and Observation
domains are only partially covered and in the process of construction.
The ancestry is particularly useful for the drug domain as it allows
browsing all drugs with a given ingredient or members of drug classes
irrespective of the country of origin, brand name or other attributes.

Users should also be aware that vocabulary updates can introduce changes
to hierarchical structures, as relationships may be added, modified, or
deprecated over time. Therefore, researchers are strongly encouraged to
version-control their vocabulary snapshot to preserve analytic
reproducibility.

\subsection{Other Relationships}\label{other-relationships}

Relationships between two different vocabularies other than mapping and
hierarchy relationships are typically of the type ``Vocabulary A -
Vocabulary B equivalent'', which is either supplied by the original
source of the vocabulary or is built de-novo. They may serve as
approximate mappings but often are less precise than the better curated
mapping relationships. High-quality equivalence relationships (such as
``Source - RxNorm equivalent'') are always duplicated by ``Maps to''
relationship.

Internal vocabulary relationships are usually supplied by the vocabulary
provider and their quality highly depends on the vocabulary. Many of
these define relationships between clinical events and can be used for
information retrieval. For example, disorders of the urethra can be
found by following the ``Finding site of'' relationship (Table~5.5):

\begin{longtable}[]{@{}
  >{\raggedright\arraybackslash}p{(\linewidth - 2\tabcolsep) * \real{0.4500}}
  >{\raggedright\arraybackslash}p{(\linewidth - 2\tabcolsep) * \real{0.5500}}@{}}
\caption{``Finding site of'' relationship of the ``Urethra,'' indicating
conditions that are situated all in this anatomical
structure.}\label{tbl-understanding-ehr-versus-claims}\tabularnewline
\toprule\noalign{}
\begin{minipage}[b]{\linewidth}\raggedright
\texttt{CONCEPT\_ID\_1}
\end{minipage} & \begin{minipage}[b]{\linewidth}\raggedright
\texttt{CONCEPT\_ID\_2}
\end{minipage} \\
\midrule\noalign{}
\endfirsthead
\toprule\noalign{}
\begin{minipage}[b]{\linewidth}\raggedright
\texttt{CONCEPT\_ID\_1}
\end{minipage} & \begin{minipage}[b]{\linewidth}\raggedright
\texttt{CONCEPT\_ID\_2}
\end{minipage} \\
\midrule\noalign{}
\endhead
\bottomrule\noalign{}
\endlastfoot
4000504 ``Urethra part'' & 36713433 ``Partial duplication of
urethra'' \\
4000504 ``Urethra part'' & 433583 ``Epispadias'' \\
4000504 ``Urethra part'' & 443533 ``Epispadias, male'' \\
4000504 ``Urethra part'' & 4005956 ``Epispadias, female'' \\
\end{longtable}

Internal relationships within a vocabulary may represent hierarchical
(for example, ``Is a'', ``RxNorm ingredient of'') connections or
non-hierarchical semantic associations such as anatomical location,
causative agent, or associated morphology. For example, within RxNorm,
relationships like ``Precise ingredient of'' and ``Has precise
ingredient'' enable navigation between drug products and their precise
ingredients.

\section{Special Situations}\label{special-situations}

\subsection{Device Coding}\label{device-coding}

Device concepts have no standardized coding scheme that could be used to
source Standard Concepts. In many source data, devices are not even
coded or contained in an external coding scheme. For this same reason,
there is currently no hierarchical system available. External standards
like GMDN and FDA's UDI database have been considered but are not yet
integrated. As a result, device concepts in OHDSI are mostly standard,
same devices have multiple standard concepts across different
vocabularies and there is no hierarchy to group terms. If you need help
with devices or want to contribute talk to the OHDSI Device Workgroup
and refer to \textbf{\{Chapter 7\}} of this book.

\subsection{Coding in Oncology}\label{coding-in-oncology}

Cancer data present unique modeling challenges due to the complexity of
diagnoses, staging, histology, metastasis, genomic features, and
treatment pathways. Please refer to the OHDSI Oncology Workgroup to
learn more about conventions.

There are several mapping principles we want to cover in this chapter:

\begin{itemize}
\tightlist
\item
  Primary cancer diagnoses are mapped to Condition domain concepts,
  mostly to SNOMED CT. ICDO-3 terms are used where SNOMED coverage is
  insufficient.
\end{itemize}

Tumor staging, grading, and metastasis details are captured using the
specialized Cancer Modifier vocabulary, which encodes structured
AJCC/UICC-based elements. Mappings in Cancer Modifier are designed to
ensure that cancer-related data: (1) preserve key clinical distinctions
(for example, metastatic vs. localized disease), (2) support
longitudinal cohort definitions (for example, new diagnosis vs.
recurrence), (3) enable harmonized analytics across registries, EHRs,
and claims data.

\begin{itemize}
\item
  Genomic abnormalities, when available, are mapped to concepts in the
  OMOP Genomic vocabulary.
\item
  Oncology-specific measurements and observations, such as tumor
  dimensions or metastasis spread, often use post-coordination
  approaches - representing the entity and its result separately - to
  align with OMOP's Measurement/Observation model.
\item
  Chemotherapy regimens are represented using the HemOnc vocabulary,
  while individual oncology drugs are mapped via RxNorm/RxNorm
  Extension.
\end{itemize}

More work is needed to refine mappings, remove duplicates, expand
support for hematologic malignancies, and integrate molecular/genomic
features.

\subsection{Coding in Psychiatry}\label{coding-in-psychiatry}

Psychiatric and neuropsychiatric data pose unique challenges for
standardization due to the complexity of symptoms, variability of
assessment tools, and evolving diagnostic frameworks. If you interested
in this research talk to the OHDSI Psychiatry Workgroup.

In the OMOP model, psychiatric assessments are primarily captured within
the Measurement and Observation domains, depending on whether the
recorded information reflects a quantitative value or a qualitative
clinical finding. Workgroup works on integrating and harmonizing
Neuropsychiatric Assessment Tools, which include standardized
psychometric scales, questionnaires, and structured interviews, into the
Vocabularies, deduplicating terms and developing a hierarchy based on
SNOMED structure to connect measurements to clinical concepts. They
consider using Thesaurus of Psychological Index Terms and Human
Phenotype Ontology (HPO), and real-world datasets (for example,
MIMIC-IV) to inform this integration.

\subsection{Coding for GIS, Exposomes and
SDOH}\label{coding-for-gis-exposomes-and-sdoh}

Environmental context, exposomes, geographic location, and social
conditions are not represented well in the OHDSI Vocabularies. If you
are interested in research, talk to the OHDSI GIS Workgroup. One of the
outputs of group is the OMOP GIS Vocabulary Package, which (22) delivers
three coordinated vocabularies: OMOP GIS for geographic units and
spatial relations, OMOP Exposome for chemicals, pollutants, toxins, and
their biological targets, and OMOP SDOH for structured
social-determinant indicators.

To accommodate these concepts, the package adds new domain identifiers
such as Geographic Feature, Environmental Feature, Socioeconomic
Feature, and Behavioral Feature. Unlike the classical OMOP domains -
essentially routing flags that direct ETL to a specific CDM table -
these new domains act solely as semantic groupers. They organize
concepts into coherent knowledge families without prescribing storage
location. Events encoded with these concepts are still recorded in the
appropriate CDM tables such as \texttt{EXTERNAL\_EXPOSURE},
\texttt{OBSERVATION}, or \texttt{MEASUREMENT} following existing
conventions.

\subsection{Microbiology and Susceptibility
Coding}\label{microbiology-and-susceptibility-coding}

There are no comprehensive conventions for microbiology coding in OHDSI.
You should refer to Themis conventions for the up-to-date guidance.
Generally, the most common scenarios involve (1) specimen collection
with a single diagnostic result (for example, Gram stain), (2) multiple
lab procedures on a single sample, (3) culture-negative findings, and
(4) one or more organisms identified and tested against antibiotics.

OMOP CDM supports this complexity through the \texttt{MEASUREMENT},
\texttt{OBSERVATION}, and \texttt{SPECIMEN} tables, with event linkages
(*\_EVENT\_ID) connecting susceptibility results to organisms and
organisms to specimens. Antibiotic susceptibility results are typically
stored as LOINC-coded MEASUREMENTs with quantitative values (for
example, MIC) and qualitative interpretations (for example, sensitive).
When coding microbiology data you should use standard concepts from
Measurement domain to populate \texttt{MEASUREMENT\_CONCEPT\_ID} (such
as susceptibility test) and Meas Value domain to populate
\texttt{VALUE\_AS\_CONCEPT\_ID} (such as detected/not detected).

\subsection{Survey Coding}\label{survey-coding}

There are no comprehensive conventions for survey coding in OHDSI. You
should refer to Survey Workgroup for the up-to-date guidance. Broadly,
surveys can be stored as Question-Answer pairs (separate concepts) or as
pre-coordinated Question-Answer (one concept). Existing survey
vocabularies, such as PPI and UK Biobank, are a mix of both. Surveys
added to the Vocabularies generally should follow broad Vocabularies
principles. For example, they should not contain negative information
and flavors of null (not reported, not specified, etc.). If they have
codes that already have standard counterparts in the Vocabularies, they
should be mapped appropriately. If you want to add your survey
instrument, please talk to the Survey Workgroup.

\subsection{Flavors of NULL}\label{flavors-of-null}

Many vocabularies contain codes that represent some form of absence of
information. For example, of the five gender concepts 8507 ``Male,''
8532 ``Female,'' 8570 ``Ambiguous,'' 8551 ``Unknown,'' and 8521
``Other'', only the first two are Standard, and the other three are
source concepts with no mapping. In the Standardized Vocabularies, there
is intentionally no distinction why a piece of information is not
available; it might be because of an active withdrawal of information by
the patient, a missing value, a value that is not defined or
standardized in some way, or the absence of a mapping record in
\texttt{CONCEPT\_RELATIONSHIP}. Any such concept is not mapped, which
corresponds to a default mapping to the Standard Concept with the
concept ID = 0.

As per Vocabularies' principles we avoid adding new flavors of NULL to
the Vocabularies and advise against using such concepts in research.

\section{Summary}\label{summary}

\begin{itemize}
\item
  All events and administrative facts are represented in the OHDSI
  Standardized Vocabularies as concepts and concept relationships.
\item
  Most of these are adopted from existing coding schemes or
  vocabularies, while others are either extended (for example, RxNorm
  Extension, OMOP Extension) or developed de novo by OHDSI Vocabulary
  Team or community to cover missing areas.
\item
  All concepts are assigned a domain, which controls where the fact
  represented by the concept is stored in the CDM.
\item
  Concepts of equivalent meaning in different vocabularies are mapped to
  one of them, which is designated the Standard Concept. The others are
  source concepts. Standard concepts (``S'') are the only concepts used
  in analytical fields.
\item
  We strive for collaborative and transparent Vocabularies with most of
  the documentation located on OHDSI Vocabularies GitHub Wiki. You can
  get involved as a community contributor or vocabulary steward. You can
  contribute simple content through templates or more complex content
  though programmatic vocabulary development.
\end{itemize}

\textbf{References}

1. Reich C, Ostropolets A, Ryan P, Rijnbeek P, Schuemie M, Davydov A, et
al.~OHDSI Standardized Vocabularies-a large-scale centralized reference
ontology for international data harmonization. J Am Med Inform Assoc.
2024 Feb 16;31(3):583--90.

2. Athena {[}Internet{]}. {[}cited 2025 May 23{]}. Available from:
https://athena.ohdsi.org/search-terms/start

3. Release planning · OHDSI/Vocabulary-v5.0 Wiki · GitHub
{[}Internet{]}. {[}cited 2025 May 23{]}. Available from:
https://github.com/OHDSI/Vocabulary-v5.0/wiki/Release-planning

4. Standardized Vocabularies · OHDSI/Vocabulary-v5.0 Wiki · GitHub
{[}Internet{]}. {[}cited 2025 May 23{]}. Available from:
https://github.com/OHDSI/Vocabulary-v5.0/wiki/Standardized-Vocabularies

5. Community contribution · OHDSI/Vocabulary-v5.0 Wiki · GitHub
{[}Internet{]}. {[}cited 2025 May 23{]}. Available from:
https://github.com/OHDSI/Vocabulary-v5.0/wiki/Community-contribution

6. Dymshyts D. Evaluating the impact of different vocabulary versions on
cohort definitions and CDM. In 2024 {[}cited 2025 May 23{]}. Available
from:
https://www.ohdsi.org/wp-content/uploads/2024/10/23-EvaluationConceptSets\_Ddymshyts\_2024\_US-Dmitry-Dymshyts.pdf

7. Amos L, Anderson D, Brody S, Ripple A, Humphreys BL. UMLS users and
uses: a current overview. Journal of the American Medical Informatics
Association. 2020 Oct 1;27(10):1606--11.

8. De Groot R, Glaser S, Kogan A, Medlock S, Alloni A, Gabetta M, et
al.~ATC-to-RxNorm mappings -- A comparison between OHDSI Standardized
Vocabularies and UMLS Metathesaurus. International Journal of Medical
Informatics. 2025 Mar;195:105777.

9. A High-Fidelity Combined ATC-Rxnorm Drug Hierarchy for Large-Scale
Observational Research. In: Studies in Health Technology and Informatics
{[}Internet{]}. IOS Press; 2024 {[}cited 2025 May 23{]}. Available from:
https://ebooks.iospress.nl/doi/10.3233/SHTI230926

10. General Structure, Download and Use · OHDSI/Vocabulary-v5.0 Wiki ·
GitHub {[}Internet{]}. {[}cited 2025 May 23{]}. Available from:
https://github.com/OHDSI/Vocabulary-v5.0/wiki/General-Structure,-Download-and-Use

11. Trofymenko M, Talapova P, Williams A. OMOP GIS Vocabulary Package
for Observational Studies in Health Care and Public Health. In.

12. GitHub {[}Internet{]}. {[}cited 2025 May 26{]}. Vocabulary
Development Process. Available from:
https://github.com/OHDSI/Vocabulary-v5.0/wiki/Vocabulary-Development-Process

13. Quality Assurance and Control · OHDSI/Vocabulary-v5.0 Wiki · GitHub
{[}Internet{]}. {[}cited 2025 May 26{]}. Available from:
https://github.com/OHDSI/Vocabulary-v5.0/wiki/Quality-assurance-and-control

14. Introduction · OHDSI/Vocabulary-v5.0 Wiki · GitHub {[}Internet{]}.
{[}cited 2025 May 26{]}. Available from:
https://github.com/OHDSI/Vocabulary-v5.0/wiki/Introduction

15. Ostropolets A. OHDSI Vocabularies landscape assessment
{[}Internet{]}. 2023. Available from:
https://ohdsiorg.sharepoint.com/:w:/s/Workgroup-CommonDataModel/EQZxds1n62JIsywmDCknwtABnSb42q7hM5PwiyblXV9zDw?e=cyvxi0

16. Releases · OHDSI/Vocabulary-v5.0 {[}Internet{]}. {[}cited 2025 May
23{]}. Available from: https://github.com/OHDSI/Vocabulary-v5.0/releases

17. Matentzoglu N, Balhoff JP, Bello SM, Bizon C, Brush M, Callahan TJ,
et al.~A Simple Standard for Sharing Ontological Mappings (SSSOM).
Database. 2022 May 25;2022:baac035.

18. OHDSI/Tantalus {[}Internet{]}. Observational Health Data Sciences
and Informatics; 2024 {[}cited 2025 May 28{]}. Available from:
https://github.com/OHDSI/Tantalus

19. Park Y, Yoon J, Zhuk A, Ostropolets A, You SC. Integrating Local
Vocabulary into OMOP CDM: A Step-by-Step Tutorial {[}Internet{]}. 2025
{[}cited 2025 May 26{]}. Available from:
http://medrxiv.org/lookup/doi/10.1101/2025.05.07.25327200

20. GitHub {[}Internet{]}. {[}cited 2025 May 23{]}. Domains. Available
from: https://github.com/OHDSI/Vocabulary-v5.0/wiki/Domains

21. Dinu V, Nadkarni P. Guidelines for the effective use of
entity--attribute--value modeling for biomedical databases.
International Journal of Medical Informatics. 2007
Nov;76(11--12):769--79.

22. OHDSI GIS WG {[}Internet{]}. 2025 {[}cited 2025 May 28{]}. Available
from: https://ohdsi.github.io/GIS/vocabulary.html

\chapter{Extract Transform Load}\label{sec-representation-etl}

TODO- write abstract

\hfill\break

\textbf{Chapter Leads}: Erica Voss

\begin{tcolorbox}[enhanced jigsaw, opacitybacktitle=0.6, coltitle=black, colback=white, bottomrule=.15mm, bottomtitle=1mm, title=\textcolor{quarto-callout-note-color}{\faInfo}\hspace{0.5em}{Note}, titlerule=0mm, rightrule=.15mm, colbacktitle=quarto-callout-note-color!10!white, colframe=quarto-callout-note-color-frame, toptitle=1mm, arc=.35mm, left=2mm, breakable, opacityback=0, leftrule=.75mm, toprule=.15mm]

This is page is currently a stub. The chapter is being written in the
OHDSI Teams directory. When the draft is complete, it will be converted
to markdown and moved to this file.

Author Resources (requires an OHDSI Teams account):

\begin{itemize}
\tightlist
\item
  \href{https://ohdsiorg.sharepoint.com/sites/Workgroup-EducationWorkingGroup/Shared\%20Documents/Forms/AllItems.aspx?id=\%2Fsites\%2FWorkgroup\%2DEducationWorkingGroup\%2FShared\%20Documents\%2FSection\%20II\%20\%2D\%20Data\%2FChapter\%206\%20\%2D\%20Extract\%20Transform\%20Load&viewid=05fec2cc\%2Dec8a\%2D4d04\%2Db565\%2Dcf1289b96f67}{Chapter
  6 Directory}
\item
  \href{https://ohdsiorg.sharepoint.com/:x:/r/sites/Workgroup-EducationWorkingGroup/_layouts/15/Doc2.aspx?action=edit&sourcedoc=\%7B1fa31e39-1c5f-4918-b878-609ebd9810b3\%7D&wdOrigin=TEAMS-WEB.teamsSdk_ns.rwc&wdExp=TEAMS-TREATMENT&wdhostclicktime=1748104477731&web=1}{Book
  Layout}
\item
  \href{https://ohdsiorg.sharepoint.com/sites/Workgroup-EducationWorkingGroup/Shared\%20Documents/Forms/AllItems.aspx?viewid=05fec2cc\%2Dec8a\%2D4d04\%2Db565\%2Dcf1289b96f67}{Education
  Working Group SharePoint Drive}
\end{itemize}

Public Resources:

\begin{itemize}
\tightlist
\item
  \href{https://ohdsi.github.io/TheBookOfOhdsi/}{Book of OHDSI, Edition
  1}
\item
  \href{https://github.com/OHDSI/TheBookOfOhdsi}{Source Code} for Book
  of OHDSI, Edition 1
\item
  \href{https://ohdsi.org/}{OHDSI Home Page}
\end{itemize}

\end{tcolorbox}

\chapter{Data Sources}\label{sec-representation-sources}

TODO- write abstract

\hfill\break

\textbf{Chapter Leads}: Melanie Philofsky

\begin{tcolorbox}[enhanced jigsaw, opacitybacktitle=0.6, coltitle=black, colback=white, bottomrule=.15mm, bottomtitle=1mm, title=\textcolor{quarto-callout-note-color}{\faInfo}\hspace{0.5em}{Note}, titlerule=0mm, rightrule=.15mm, colbacktitle=quarto-callout-note-color!10!white, colframe=quarto-callout-note-color-frame, toptitle=1mm, arc=.35mm, left=2mm, breakable, opacityback=0, leftrule=.75mm, toprule=.15mm]

This is page is currently a stub. The chapter is being written in the
OHDSI Teams directory. When the draft is complete, it will be converted
to markdown and moved to this file.

Author Resources (requires an OHDSI Teams account):

\begin{itemize}
\tightlist
\item
  \href{https://ohdsiorg.sharepoint.com/sites/Workgroup-EducationWorkingGroup/Shared\%20Documents/Forms/AllItems.aspx?id=\%2Fsites\%2FWorkgroup\%2DEducationWorkingGroup\%2FShared\%20Documents\%2FSection\%20II\%20\%2D\%20Data\%2FChapter\%207\%20\%2D\%20Data\%20Sources&viewid=05fec2cc\%2Dec8a\%2D4d04\%2Db565\%2Dcf1289b96f67}{Chapter
  7 Directory}
\item
  \href{https://ohdsiorg.sharepoint.com/:x:/r/sites/Workgroup-EducationWorkingGroup/_layouts/15/Doc2.aspx?action=edit&sourcedoc=\%7B1fa31e39-1c5f-4918-b878-609ebd9810b3\%7D&wdOrigin=TEAMS-WEB.teamsSdk_ns.rwc&wdExp=TEAMS-TREATMENT&wdhostclicktime=1748104477731&web=1}{Book
  Layout}
\item
  \href{https://ohdsiorg.sharepoint.com/sites/Workgroup-EducationWorkingGroup/Shared\%20Documents/Forms/AllItems.aspx?viewid=05fec2cc\%2Dec8a\%2D4d04\%2Db565\%2Dcf1289b96f67}{Education
  Working Group SharePoint Drive}
\end{itemize}

Public Resources:

\begin{itemize}
\tightlist
\item
  \href{https://ohdsi.github.io/TheBookOfOhdsi/}{Book of OHDSI, Edition
  1}
\item
  \href{https://github.com/OHDSI/TheBookOfOhdsi}{Source Code} for Book
  of OHDSI, Edition 1
\item
  \href{https://ohdsi.org/}{OHDSI Home Page}
\end{itemize}

\end{tcolorbox}

\part{Data Analytics}

\chapter{Data Analytics Use Cases}\label{sec-analytics-cases}

TODO- write abstract

\hfill\break

\textbf{Chapter Leads}: Rakesh Babu

\begin{tcolorbox}[enhanced jigsaw, opacitybacktitle=0.6, coltitle=black, colback=white, bottomrule=.15mm, bottomtitle=1mm, title=\textcolor{quarto-callout-note-color}{\faInfo}\hspace{0.5em}{Note}, titlerule=0mm, rightrule=.15mm, colbacktitle=quarto-callout-note-color!10!white, colframe=quarto-callout-note-color-frame, toptitle=1mm, arc=.35mm, left=2mm, breakable, opacityback=0, leftrule=.75mm, toprule=.15mm]

This is page is currently a stub. The chapter is being written in the
OHDSI Teams directory. When the draft is complete, it will be converted
to markdown and moved to this file.

Author Resources (requires an OHDSI Teams account):

\begin{itemize}
\tightlist
\item
  \href{https://ohdsiorg.sharepoint.com/sites/Workgroup-EducationWorkingGroup/Shared\%20Documents/Forms/AllItems.aspx?id=\%2Fsites\%2FWorkgroup\%2DEducationWorkingGroup\%2FShared\%20Documents\%2FSection\%20III\%20\%2D\%20Analytics\%2FChapter\%208\%20Introduction&viewid=05fec2cc\%2Dec8a\%2D4d04\%2Db565\%2Dcf1289b96f67}{Chapter
  Directory}
\item
  \href{https://ohdsiorg.sharepoint.com/:x:/r/sites/Workgroup-EducationWorkingGroup/_layouts/15/Doc2.aspx?action=edit&sourcedoc=\%7B1fa31e39-1c5f-4918-b878-609ebd9810b3\%7D&wdOrigin=TEAMS-WEB.teamsSdk_ns.rwc&wdExp=TEAMS-TREATMENT&wdhostclicktime=1748104477731&web=1}{Book
  Layout}
\item
  \href{https://ohdsiorg.sharepoint.com/sites/Workgroup-EducationWorkingGroup/Shared\%20Documents/Forms/AllItems.aspx?viewid=05fec2cc\%2Dec8a\%2D4d04\%2Db565\%2Dcf1289b96f67}{Education
  Working Group SharePoint Drive}
\end{itemize}

Public Resources:

\begin{itemize}
\tightlist
\item
  \href{https://ohdsi.github.io/TheBookOfOhdsi/}{Book of OHDSI, Edition
  1}
\item
  \href{https://github.com/OHDSI/TheBookOfOhdsi}{Source Code} for Book
  of OHDSI, Edition 1
\item
  \href{https://ohdsi.org/}{OHDSI Home Page}
\end{itemize}

\end{tcolorbox}

\chapter{OHDSI Analytics Tools}\label{sec-analytics-tools}

TODO- write abstract

\hfill\break

\textbf{Chapter Leads}: Anthony Sena

\begin{tcolorbox}[enhanced jigsaw, opacitybacktitle=0.6, coltitle=black, colback=white, bottomrule=.15mm, bottomtitle=1mm, title=\textcolor{quarto-callout-note-color}{\faInfo}\hspace{0.5em}{Note}, titlerule=0mm, rightrule=.15mm, colbacktitle=quarto-callout-note-color!10!white, colframe=quarto-callout-note-color-frame, toptitle=1mm, arc=.35mm, left=2mm, breakable, opacityback=0, leftrule=.75mm, toprule=.15mm]

This is page is currently a stub. The chapter is being written in the
OHDSI Teams directory. When the draft is complete, it will be converted
to markdown and moved to this file.

Author Resources (requires an OHDSI Teams account):

\begin{itemize}
\tightlist
\item
  \href{https://ohdsiorg.sharepoint.com/sites/Workgroup-EducationWorkingGroup/Shared\%20Documents/Forms/AllItems.aspx?id=\%2Fsites\%2FWorkgroup\%2DEducationWorkingGroup\%2FShared\%20Documents\%2FSection\%20III\%20\%2D\%20Analytics\%2FChapter\%209\%20OHDSI\%20Analytical\%20Tools&viewid=05fec2cc\%2Dec8a\%2D4d04\%2Db565\%2Dcf1289b96f67}{Chapter
  Directory}
\item
  \href{https://ohdsiorg.sharepoint.com/:x:/r/sites/Workgroup-EducationWorkingGroup/_layouts/15/Doc2.aspx?action=edit&sourcedoc=\%7B1fa31e39-1c5f-4918-b878-609ebd9810b3\%7D&wdOrigin=TEAMS-WEB.teamsSdk_ns.rwc&wdExp=TEAMS-TREATMENT&wdhostclicktime=1748104477731&web=1}{Book
  Layout}
\item
  \href{https://ohdsiorg.sharepoint.com/sites/Workgroup-EducationWorkingGroup/Shared\%20Documents/Forms/AllItems.aspx?viewid=05fec2cc\%2Dec8a\%2D4d04\%2Db565\%2Dcf1289b96f67}{Education
  Working Group SharePoint Drive}
\end{itemize}

Public Resources:

\begin{itemize}
\tightlist
\item
  \href{https://ohdsi.github.io/TheBookOfOhdsi/}{Book of OHDSI, Edition
  1}
\item
  \href{https://github.com/OHDSI/TheBookOfOhdsi}{Source Code} for Book
  of OHDSI, Edition 1
\item
  \href{https://ohdsi.org/}{OHDSI Home Page}
\end{itemize}

\end{tcolorbox}

\chapter{Defining Cohorts}\label{sec-analytics-cohorts}

TODO- write abstract

\hfill\break

\textbf{Chapter Leads}: Azza Shoaibi

\begin{tcolorbox}[enhanced jigsaw, opacitybacktitle=0.6, coltitle=black, colback=white, bottomrule=.15mm, bottomtitle=1mm, title=\textcolor{quarto-callout-note-color}{\faInfo}\hspace{0.5em}{Note}, titlerule=0mm, rightrule=.15mm, colbacktitle=quarto-callout-note-color!10!white, colframe=quarto-callout-note-color-frame, toptitle=1mm, arc=.35mm, left=2mm, breakable, opacityback=0, leftrule=.75mm, toprule=.15mm]

This is page is currently a stub. The chapter is being written in the
OHDSI Teams directory. When the draft is complete, it will be converted
to markdown and moved to this file.

Author Resources (requires an OHDSI Teams account):

\begin{itemize}
\tightlist
\item
  \href{https://ohdsiorg.sharepoint.com/sites/Workgroup-EducationWorkingGroup/Shared\%20Documents/Forms/AllItems.aspx?id=\%2Fsites\%2FWorkgroup\%2DEducationWorkingGroup\%2FShared\%20Documents\%2FSection\%20III\%20\%2D\%20Analytics\%2FChapter\%2011\%20Defining\%20Cohorts&viewid=05fec2cc\%2Dec8a\%2D4d04\%2Db565\%2Dcf1289b96f67}{Chapter
  Directory}
\item
  \href{https://ohdsiorg.sharepoint.com/:x:/r/sites/Workgroup-EducationWorkingGroup/_layouts/15/Doc2.aspx?action=edit&sourcedoc=\%7B1fa31e39-1c5f-4918-b878-609ebd9810b3\%7D&wdOrigin=TEAMS-WEB.teamsSdk_ns.rwc&wdExp=TEAMS-TREATMENT&wdhostclicktime=1748104477731&web=1}{Book
  Layout}
\item
  \href{https://ohdsiorg.sharepoint.com/sites/Workgroup-EducationWorkingGroup/Shared\%20Documents/Forms/AllItems.aspx?viewid=05fec2cc\%2Dec8a\%2D4d04\%2Db565\%2Dcf1289b96f67}{Education
  Working Group SharePoint Drive}
\end{itemize}

Public Resources:

\begin{itemize}
\tightlist
\item
  \href{https://ohdsi.github.io/TheBookOfOhdsi/}{Book of OHDSI, Edition
  1}
\item
  \href{https://github.com/OHDSI/TheBookOfOhdsi}{Source Code} for Book
  of OHDSI, Edition 1
\item
  \href{https://ohdsi.org/}{OHDSI Home Page}
\end{itemize}

\end{tcolorbox}

\chapter{Characterization}\label{sec-analytics-characterization}

TODO- write abstract

\hfill\break

\textbf{Chapter Leads}: Gowtham Rao

\begin{tcolorbox}[enhanced jigsaw, opacitybacktitle=0.6, coltitle=black, colback=white, bottomrule=.15mm, bottomtitle=1mm, title=\textcolor{quarto-callout-note-color}{\faInfo}\hspace{0.5em}{Note}, titlerule=0mm, rightrule=.15mm, colbacktitle=quarto-callout-note-color!10!white, colframe=quarto-callout-note-color-frame, toptitle=1mm, arc=.35mm, left=2mm, breakable, opacityback=0, leftrule=.75mm, toprule=.15mm]

This is page is currently a stub. The chapter is being written in the
OHDSI Teams directory. When the draft is complete, it will be converted
to markdown and moved to this file.

Author Resources (requires an OHDSI Teams account):

\begin{itemize}
\tightlist
\item
  \href{https://ohdsiorg.sharepoint.com/sites/Workgroup-EducationWorkingGroup/Shared\%20Documents/Forms/AllItems.aspx?id=\%2Fsites\%2FWorkgroup\%2DEducationWorkingGroup\%2FShared\%20Documents\%2FSection\%20III\%20\%2D\%20Analytics\%2FChapter\%2012\%20Characterization&viewid=05fec2cc\%2Dec8a\%2D4d04\%2Db565\%2Dcf1289b96f67}{Chapter
  Directory}
\item
  \href{https://ohdsiorg.sharepoint.com/:x:/r/sites/Workgroup-EducationWorkingGroup/_layouts/15/Doc2.aspx?action=edit&sourcedoc=\%7B1fa31e39-1c5f-4918-b878-609ebd9810b3\%7D&wdOrigin=TEAMS-WEB.teamsSdk_ns.rwc&wdExp=TEAMS-TREATMENT&wdhostclicktime=1748104477731&web=1}{Book
  Layout}
\item
  \href{https://ohdsiorg.sharepoint.com/sites/Workgroup-EducationWorkingGroup/Shared\%20Documents/Forms/AllItems.aspx?viewid=05fec2cc\%2Dec8a\%2D4d04\%2Db565\%2Dcf1289b96f67}{Education
  Working Group SharePoint Drive}
\end{itemize}

Public Resources:

\begin{itemize}
\tightlist
\item
  \href{https://ohdsi.github.io/TheBookOfOhdsi/}{Book of OHDSI, Edition
  1}
\item
  \href{https://github.com/OHDSI/TheBookOfOhdsi}{Source Code} for Book
  of OHDSI, Edition 1
\item
  \href{https://ohdsi.org/}{OHDSI Home Page}
\end{itemize}

\end{tcolorbox}

\chapter{Population-Level Estimation}\label{sec-analytics-population}

TODO- write abstract

\hfill\break

\textbf{Chapter Leads}: Martijn Schuemie

\begin{tcolorbox}[enhanced jigsaw, opacitybacktitle=0.6, coltitle=black, colback=white, bottomrule=.15mm, bottomtitle=1mm, title=\textcolor{quarto-callout-note-color}{\faInfo}\hspace{0.5em}{Note}, titlerule=0mm, rightrule=.15mm, colbacktitle=quarto-callout-note-color!10!white, colframe=quarto-callout-note-color-frame, toptitle=1mm, arc=.35mm, left=2mm, breakable, opacityback=0, leftrule=.75mm, toprule=.15mm]

This is page is currently a stub. The chapter is being written in the
OHDSI Teams directory. When the draft is complete, it will be converted
to markdown and moved to this file.

Author Resources (requires an OHDSI Teams account):

\begin{itemize}
\tightlist
\item
  \href{https://ohdsiorg.sharepoint.com/sites/Workgroup-EducationWorkingGroup/Shared\%20Documents/Forms/AllItems.aspx?id=\%2Fsites\%2FWorkgroup\%2DEducationWorkingGroup\%2FShared\%20Documents\%2FSection\%20III\%20\%2D\%20Analytics\%2FChapter\%2013\%20PLE&viewid=05fec2cc\%2Dec8a\%2D4d04\%2Db565\%2Dcf1289b96f67}{Chapter
  Directory}
\item
  \href{https://ohdsiorg.sharepoint.com/:x:/r/sites/Workgroup-EducationWorkingGroup/_layouts/15/Doc2.aspx?action=edit&sourcedoc=\%7B1fa31e39-1c5f-4918-b878-609ebd9810b3\%7D&wdOrigin=TEAMS-WEB.teamsSdk_ns.rwc&wdExp=TEAMS-TREATMENT&wdhostclicktime=1748104477731&web=1}{Book
  Layout}
\item
  \href{https://ohdsiorg.sharepoint.com/sites/Workgroup-EducationWorkingGroup/Shared\%20Documents/Forms/AllItems.aspx?viewid=05fec2cc\%2Dec8a\%2D4d04\%2Db565\%2Dcf1289b96f67}{Education
  Working Group SharePoint Drive}
\end{itemize}

Public Resources:

\begin{itemize}
\tightlist
\item
  \href{https://ohdsi.github.io/TheBookOfOhdsi/}{Book of OHDSI, Edition
  1}
\item
  \href{https://github.com/OHDSI/TheBookOfOhdsi}{Source Code} for Book
  of OHDSI, Edition 1
\item
  \href{https://ohdsi.org/}{OHDSI Home Page}
\end{itemize}

\end{tcolorbox}

\chapter{Patient-Level Prediction}\label{sec-analytics-patient}

TODO- write abstract

\hfill\break

\textbf{Chapter Leads}: Ross Williams

\begin{tcolorbox}[enhanced jigsaw, opacitybacktitle=0.6, coltitle=black, colback=white, bottomrule=.15mm, bottomtitle=1mm, title=\textcolor{quarto-callout-note-color}{\faInfo}\hspace{0.5em}{Note}, titlerule=0mm, rightrule=.15mm, colbacktitle=quarto-callout-note-color!10!white, colframe=quarto-callout-note-color-frame, toptitle=1mm, arc=.35mm, left=2mm, breakable, opacityback=0, leftrule=.75mm, toprule=.15mm]

This is page is currently a stub. The chapter is being written in the
OHDSI Teams directory. When the draft is complete, it will be converted
to markdown and moved to this file.

Author Resources (requires an OHDSI Teams account):

\begin{itemize}
\tightlist
\item
  \href{https://ohdsiorg.sharepoint.com/sites/Workgroup-EducationWorkingGroup/Shared\%20Documents/Forms/AllItems.aspx?id=\%2Fsites\%2FWorkgroup\%2DEducationWorkingGroup\%2FShared\%20Documents\%2FSection\%20III\%20\%2D\%20Analytics\%2FChapter\%2014\%20PLP&viewid=05fec2cc\%2Dec8a\%2D4d04\%2Db565\%2Dcf1289b96f67}{Chapter
  Directory}
\item
  \href{https://ohdsiorg.sharepoint.com/:x:/r/sites/Workgroup-EducationWorkingGroup/_layouts/15/Doc2.aspx?action=edit&sourcedoc=\%7B1fa31e39-1c5f-4918-b878-609ebd9810b3\%7D&wdOrigin=TEAMS-WEB.teamsSdk_ns.rwc&wdExp=TEAMS-TREATMENT&wdhostclicktime=1748104477731&web=1}{Book
  Layout}
\item
  \href{https://ohdsiorg.sharepoint.com/sites/Workgroup-EducationWorkingGroup/Shared\%20Documents/Forms/AllItems.aspx?viewid=05fec2cc\%2Dec8a\%2D4d04\%2Db565\%2Dcf1289b96f67}{Education
  Working Group SharePoint Drive}
\end{itemize}

Public Resources:

\begin{itemize}
\tightlist
\item
  \href{https://ohdsi.github.io/TheBookOfOhdsi/}{Book of OHDSI, Edition
  1}
\item
  \href{https://github.com/OHDSI/TheBookOfOhdsi}{Source Code} for Book
  of OHDSI, Edition 1
\item
  \href{https://ohdsi.org/}{OHDSI Home Page}
\end{itemize}

\end{tcolorbox}

\part{Evidence Quality}

\chapter{Evidence Quality}\label{sec-quality-evidence}

TODO- write abstract

\hfill\break

\textbf{Chapter Leads}: Patrick Ryan

\begin{tcolorbox}[enhanced jigsaw, opacitybacktitle=0.6, coltitle=black, colback=white, bottomrule=.15mm, bottomtitle=1mm, title=\textcolor{quarto-callout-note-color}{\faInfo}\hspace{0.5em}{Note}, titlerule=0mm, rightrule=.15mm, colbacktitle=quarto-callout-note-color!10!white, colframe=quarto-callout-note-color-frame, toptitle=1mm, arc=.35mm, left=2mm, breakable, opacityback=0, leftrule=.75mm, toprule=.15mm]

This is page is currently a stub. The chapter is being written in the
OHDSI Teams directory. When the draft is complete, it will be converted
to markdown and moved to this file.

Author Resources (requires an OHDSI Teams account):

\begin{itemize}
\tightlist
\item
  \href{https://ohdsiorg.sharepoint.com/sites/Workgroup-EducationWorkingGroup/Shared\%20Documents/Forms/AllItems.aspx?id=\%2Fsites\%2FWorkgroup\%2DEducationWorkingGroup\%2FShared\%20Documents\%2FSection\%20IV\%20\%2D\%20Quality\%2FChapter\%2015\%20\%2D\%20Evidence\%20Quality&viewid=05fec2cc\%2Dec8a\%2D4d04\%2Db565\%2Dcf1289b96f67}{Chapter
  Directory}
\item
  \href{https://ohdsiorg.sharepoint.com/:x:/r/sites/Workgroup-EducationWorkingGroup/_layouts/15/Doc2.aspx?action=edit&sourcedoc=\%7B1fa31e39-1c5f-4918-b878-609ebd9810b3\%7D&wdOrigin=TEAMS-WEB.teamsSdk_ns.rwc&wdExp=TEAMS-TREATMENT&wdhostclicktime=1748104477731&web=1}{Book
  Layout}
\item
  \href{https://ohdsiorg.sharepoint.com/sites/Workgroup-EducationWorkingGroup/Shared\%20Documents/Forms/AllItems.aspx?viewid=05fec2cc\%2Dec8a\%2D4d04\%2Db565\%2Dcf1289b96f67}{Education
  Working Group SharePoint Drive}
\end{itemize}

Public Resources:

\begin{itemize}
\tightlist
\item
  \href{https://ohdsi.github.io/TheBookOfOhdsi/}{Book of OHDSI, Edition
  1}
\item
  \href{https://github.com/OHDSI/TheBookOfOhdsi}{Source Code} for Book
  of OHDSI, Edition 1
\item
  \href{https://ohdsi.org/}{OHDSI Home Page}
\end{itemize}

\end{tcolorbox}

\chapter{Data Quality}\label{sec-quality-data}

TODO- write abstract

\hfill\break

\textbf{Chapter Leads}: Clair Blacketer

\begin{tcolorbox}[enhanced jigsaw, opacitybacktitle=0.6, coltitle=black, colback=white, bottomrule=.15mm, bottomtitle=1mm, title=\textcolor{quarto-callout-note-color}{\faInfo}\hspace{0.5em}{Note}, titlerule=0mm, rightrule=.15mm, colbacktitle=quarto-callout-note-color!10!white, colframe=quarto-callout-note-color-frame, toptitle=1mm, arc=.35mm, left=2mm, breakable, opacityback=0, leftrule=.75mm, toprule=.15mm]

This is page is currently a stub. The chapter is being written in the
OHDSI Teams directory. When the draft is complete, it will be converted
to markdown and moved to this file.

Author Resources (requires an OHDSI Teams account):

\begin{itemize}
\tightlist
\item
  \href{https://ohdsiorg.sharepoint.com/sites/Workgroup-EducationWorkingGroup/Shared\%20Documents/Forms/AllItems.aspx?id=\%2Fsites\%2FWorkgroup\%2DEducationWorkingGroup\%2FShared\%20Documents\%2FSection\%20IV\%20\%2D\%20Quality\%2FChapter\%2016\%20\%2D\%20Data\%20Quality&viewid=05fec2cc\%2Dec8a\%2D4d04\%2Db565\%2Dcf1289b96f67}{Chapter
  Directory}
\item
  \href{https://ohdsiorg.sharepoint.com/:x:/r/sites/Workgroup-EducationWorkingGroup/_layouts/15/Doc2.aspx?action=edit&sourcedoc=\%7B1fa31e39-1c5f-4918-b878-609ebd9810b3\%7D&wdOrigin=TEAMS-WEB.teamsSdk_ns.rwc&wdExp=TEAMS-TREATMENT&wdhostclicktime=1748104477731&web=1}{Book
  Layout}
\item
  \href{https://ohdsiorg.sharepoint.com/sites/Workgroup-EducationWorkingGroup/Shared\%20Documents/Forms/AllItems.aspx?viewid=05fec2cc\%2Dec8a\%2D4d04\%2Db565\%2Dcf1289b96f67}{Education
  Working Group SharePoint Drive}
\end{itemize}

Public Resources:

\begin{itemize}
\tightlist
\item
  \href{https://ohdsi.github.io/TheBookOfOhdsi/}{Book of OHDSI, Edition
  1}
\item
  \href{https://github.com/OHDSI/TheBookOfOhdsi}{Source Code} for Book
  of OHDSI, Edition 1
\item
  \href{https://ohdsi.org/}{OHDSI Home Page}
\end{itemize}

\end{tcolorbox}

\chapter{Clinical Validity}\label{sec-quality-clinical}

TODO- write abstract

\hfill\break

\textbf{Chapter Leads}: Joel Swerdel

\begin{tcolorbox}[enhanced jigsaw, opacitybacktitle=0.6, coltitle=black, colback=white, bottomrule=.15mm, bottomtitle=1mm, title=\textcolor{quarto-callout-note-color}{\faInfo}\hspace{0.5em}{Note}, titlerule=0mm, rightrule=.15mm, colbacktitle=quarto-callout-note-color!10!white, colframe=quarto-callout-note-color-frame, toptitle=1mm, arc=.35mm, left=2mm, breakable, opacityback=0, leftrule=.75mm, toprule=.15mm]

This is page is currently a stub. The chapter is being written in the
OHDSI Teams directory. When the draft is complete, it will be converted
to markdown and moved to this file.

Author Resources (requires an OHDSI Teams account):

\begin{itemize}
\tightlist
\item
  \href{https://ohdsiorg.sharepoint.com/sites/Workgroup-EducationWorkingGroup/Shared\%20Documents/Forms/AllItems.aspx?id=\%2Fsites\%2FWorkgroup\%2DEducationWorkingGroup\%2FShared\%20Documents\%2FSection\%20IV\%20\%2D\%20Quality\%2FChapter\%2017\%20\%2D\%20Cohort\%20Validity&viewid=05fec2cc\%2Dec8a\%2D4d04\%2Db565\%2Dcf1289b96f67}{Chapter
  Directory}
\item
  \href{https://ohdsiorg.sharepoint.com/:x:/r/sites/Workgroup-EducationWorkingGroup/_layouts/15/Doc2.aspx?action=edit&sourcedoc=\%7B1fa31e39-1c5f-4918-b878-609ebd9810b3\%7D&wdOrigin=TEAMS-WEB.teamsSdk_ns.rwc&wdExp=TEAMS-TREATMENT&wdhostclicktime=1748104477731&web=1}{Book
  Layout}
\item
  \href{https://ohdsiorg.sharepoint.com/sites/Workgroup-EducationWorkingGroup/Shared\%20Documents/Forms/AllItems.aspx?viewid=05fec2cc\%2Dec8a\%2D4d04\%2Db565\%2Dcf1289b96f67}{Education
  Working Group SharePoint Drive}
\end{itemize}

Public Resources:

\begin{itemize}
\tightlist
\item
  \href{https://ohdsi.github.io/TheBookOfOhdsi/}{Book of OHDSI, Edition
  1}
\item
  \href{https://github.com/OHDSI/TheBookOfOhdsi}{Source Code} for Book
  of OHDSI, Edition 1
\item
  \href{https://ohdsi.org/}{OHDSI Home Page}
\end{itemize}

\end{tcolorbox}

\begin{quote}
The likelihood of transforming matter into energy is something akin to
shooting birds in the dark in a country where there are only a few
birds. \emph{Einstein, 1935}
\end{quote}

The vision of OHDSI is ``A world in which observational research
produces a comprehensive understanding of health and disease.''
Retrospective designs provide a vehicle for research using existing data
but can be riddled with threats to various aspects of validity as
discussed in Chapter14. It is not easy to isolate clinical validity from
quality of data and statistical methodology, but here we will focus on
three aspects in terms of clinical validity: Characteristics of health
care databases, Cohort validation, and Generalizability of the evidence.
Let's go back to the example of population-level estimation (Chapter
12). We tried to answer the question ``Do ACE inhibitors cause
angioedema compared to thiazide or thiazide-like diuretics?'' In that
example, we demonstrated that ACE inhibitors caused more angioedema than
thiazide or thiazide-like diuretics. This chapter is dedicated to answer
the question: ``To what extent does the analysis conducted match the
clinical intention?''

\section{Characteristics of Health Care
Databases}\label{characteristics-of-health-care-databases}

It is possible that what we found is the relationship between
\textbf{prescription} of ACE inhibitor and angioedema rather than the
relationship between \textbf{use} of ACE inhibitor and angioedema. We've
already discussed data quality in the previous chapter (15). The quality
of the converted database into the Common Data Model (CDM) cannot exceed
the original database. Here we are addressing the characteristics of
most healthcare utilization databases. Many databases used in OHDSI
originated from administrative claims or electronic health records
(EHR). Claims and EHR have different data capture processes, neither of
which has research as a primary intention. Data elements from claims
records are captured for the purpose of reimbursement, financial
transactions between clinicians and payers whereby services provided to
patients by providers are sufficiently justified to enable agreement on
payments by the responsible parties. Data elements in EHR records are
captured to support clinical care and administrative operations, and
they commonly only reflect the information that providers within a given
health system feel are necessary to document the current service and
provide necessary context for anticipated follow-up care within their
health system. They may not represent a patient's complete medical
history and may not integrate data from across health systems.

To generate reliable evidence from observational data, it is useful for
a researcher to understand the journey that the data undergoes from the
moment that a patient seeks care through the moment that the data
reflecting that care are used in an analysis. As an example, ``drug
exposure'' can be inferred from various sources of observational data,
including prescriptions written by clinicians, pharmacy dispensing
records, hospital procedural administrations, or patient self-reported
medication history. The source of data can impact our level of
confidence in the inference we draw about which patients did or did not
use the drug, as well as when and for how long. The data capture process
can result in under-estimation of exposure, such as if free samples or
over-the counter drugs are not recorded, or over-estimation of exposure,
such as if a patient doesn't fill the prescription written or doesn't
adherently consume the prescription dispensed. Understanding the
potential biases in exposure and outcome ascertainment, and more ideally
quantifying and adjusting for these measurement errors, can improve our
confidence in the validity of the evidence we draw from the data we have
available.

\section{Cohort Validation}\label{cohort-validation}

G. Hripcsak and Albers described that ``a phenotype is a specification
of an observable, potentially changing state of an organism, as
distinguished from the genotype, which is derived from an organism's
genetic makeup''. (1) The term phenotype can be applied to patient
characteristics inferred from electronic health record (EHR) data.
Researchers have been carrying out EHR phenotyping since the beginning
of informatics, from both structured data and narrative data. The goal
is to draw conclusions about a target concept based on raw EHR data,
claims data, or other clinically relevant data. Phenotype algorithms --
i.e., algorithms that identify or characterize phenotypes -- may be
generated by domain exerts and knowledge engineers, including recent
research in knowledge engineering or through diverse forms of machine
learning\ldots to generate novel representations of the data.''

This description highlights several attributes useful to reinforce when
considering clinical validity: 1) it makes it clear that we are talking
about something that is observable (and therefore possible to be
captured in our observational data); 2) it includes the notion of time
in the phenotype specification (since a state of a person can change);
3) it draws a distinction between the phenotype as the desired intent
vs.~the phenotype algorithm, which is the implementation of the desired
intent.

OHDSI has adopted the term ``cohort'' to define the set of persons
satisfying one or more inclusion criteria for a duration of time. A
``cohort definition'' represents the logic necessary to instantiate a
cohort against an observational database. In this regard, the cohort
definition (or phenotype algorithm) is used to produce a cohort, which
is intended to represent the phenotype, being the persons who belong to
the observable clinical state of interest.

Most types of observational analyses, including clinical
characterization, population-level effect estimation, and patient-level
prediction, require one or more cohorts to be established as part of the
study process. To evaluate the validity of the evidence produced by
these analyses, one must consider this question for each cohort: to what
extent do the persons identified in the cohort based on the cohort
definition and the available observational data accurately reflect the
persons who truly belong to the phenotype?

To return to the population-level estimation example (Chapter 12) ``Do
ACE inhibitors cause angioedema compared to thiazide or thiazide-like
diuretics?'', we must define three cohorts: a target cohort of persons
who are new users of ACE inhibitors, a comparator cohort of persons who
are new users of thiazide diuretics, and an outcome cohort of persons
who develop angioedema. How confident are we that all use of ACE
inhibitors or thiazide diuretics is completely captured, such that ``new
users'' can be identified by the first observed exposure, without
concern of prior (but unobserved) use? Can we comfortably infer that
persons who have a drug exposure record for ACE inhibitors were in fact
exposed to the drug, and those without a drug exposure were indeed
unexposed? Is there uncertainty in defining the duration of time that a
person is classified in the state of ``ACE inhibitor use,'' either when
inferring cohort entry at the time the drug was started or cohort exit
when the drug was discontinued? Have persons with a condition occurrence
record of ``Angioedema'' actually experienced rapid swelling beneath the
skin, differentiated from other types of dermatologic allergic
reactions? What proportion of patients who developed angioedema received
medical attention that would give rise to the observational data used to
identify these clinical cases based on the cohort definition? How well
can the angioedema events which are potentially drug-induced be
disambiguated from the events known to be caused by other agents, such
as food allergy or viral infection? Is disease onset sufficiently well
captured that we have confidence in drawing a temporal association
between exposure status and outcome incidence? Answering these types of
questions is at the heart of clinical validity.

In this chapter, we will discuss the methods for validating cohort
definitions. We first describe the metrics used to measure the validity
of a cohort definition. Next, we describe two methods to estimate these
metrics: 1) clinical adjudication through source record verification,
and 2) PheValuator, a semi-automated method using diagnostic predictive
modeling.

\subsection{Cohort Evaluation Metrics}\label{cohort-evaluation-metrics}

Once the cohort definition for the study has been determined, the
validity of the definition can be evaluated. A common approach to assess
validity is by comparing some or all persons in a defined cohort to a
reference `gold standard' and expressing the results in a confusion
matrix, a two-by-two contingency table that stratifies persons according
to their gold standard classification and qualification within the
cohort definition. Figure~\ref{fig-quality-clinical-010-confusion} shows
the elements of the confusion matrix.

\begin{figure}[H]

\centering{

\pandocbounded{\includegraphics[keepaspectratio]{quality/images/clinical/fig-quality-clinical-010-confusion.png}}

}

\caption{\label{fig-quality-clinical-010-confusion}Confusion matrix.}

\end{figure}%

The true and false results from the cohort definition are determined by
applying the definition to a group of persons. Those included in the
definition are considered positive for the health condition and are
labeled ``True.'' Those persons not included in the cohort definition
are considered negative for the health condition and are labeled
``False''. While the absolute truth of a person's heath state considered
in the cohort definition is very difficult to determine, there are
multiple methods to establish a reference gold standard, two of which
will be described later in the chapter. Regardless of the method used,
the labeling of these persons is the same as described for the cohort
definition.

In addition to errors in the binary indication of phenotype designation,
the timing of the health condition may also be incorrect. For example,
while the cohort definition may correctly label a person as belonging to
a phenotype, the definition may incorrectly specify the date and time
when a person without the condition became a person with the condition.
This error would add bias to studies using survival analysis results,
e.g., hazard ratios, as an effect measure.

The next step in the process is to assess the concordance of the gold
standard with the cohort definition. Those persons that are labeled by
both the gold standard method and the cohort definition as ``True'' are
called ``True Positives.'' Those persons that are labeled by the gold
standard method as ``False'' and by the cohort definition as ``True''
are called ``False Positives,'' i.e., the cohort definition
misclassified these persons as having the condition when they do not.
Those persons that are labeled by both the gold standard method and the
cohort definition as ``False'' are called ``True Negatives.'' Those
persons that are labeled by the gold standard method as ``True'' and by
the cohort definition as ``False'' are called ``False Negatives,'' i.e.,
the cohort definition incorrectly classified these persons as not having
the condition, when it fact they do belong to the phenotype. Using the
counts from the four cells in the confusion matrix, we can quantify the
accuracy of the cohort definition in classifying phenotype status in a
group of persons. There are standard performance metrics for measuring
cohort definition performance:

\begin{enumerate}
\def\labelenumi{\arabic{enumi}.}
\tightlist
\item
  \textbf{Sensitivity of the cohort definition} -- what proportion of
  the persons who truly belong to the phenotype in the population were
  correctly identified to have the health outcome based on the cohort
  definition? This is determined by the following formula:
\end{enumerate}

\begin{itemize}
\tightlist
\item
  Sensitivity = True Positives / (True Positives + False Negatives)
\end{itemize}

\begin{enumerate}
\def\labelenumi{\arabic{enumi}.}
\tightlist
\item
  \textbf{Specificity of the cohort definition} -- what proportion of
  the persons who do not belong to the phenotype in the population were
  correctly identified to not have the health outcome based on the
  cohort definition? This is determined by the following formula:
\end{enumerate}

\begin{itemize}
\tightlist
\item
  Specificity = True Negatives / (True Negatives + False Positives)
\end{itemize}

\begin{enumerate}
\def\labelenumi{\arabic{enumi}.}
\tightlist
\item
  \textbf{Positive predictive value (PPV) of the cohort definition} --
  what proportion of the persons identified by the cohort definition to
  have the health condition actually belong to the phenotype? This is
  determined by the following formula:
\end{enumerate}

\begin{itemize}
\tightlist
\item
  PPV = True Positives / (True Positives + False Positives)
\end{itemize}

\begin{enumerate}
\def\labelenumi{\arabic{enumi}.}
\tightlist
\item
  \textbf{Negative predictive value (NPV) of the cohort definition} --
  what proportion of the persons identified by the cohort definition to
  not have the health condition actually did not belong to the
  phenotype? This is determined by the following formula:
\end{enumerate}

\begin{itemize}
\tightlist
\item
  NPV = True Negatives / (True Negatives + False Negatives)
\end{itemize}

Perfect scores for these measures are 100\%. Due to the nature of
observational data, perfect scores are usually far from the norm. Rubbo
et al.~reviewed studies validating cohort definitions for myocardial
infarction. (2) Of the 33 studies they examined, only one cohort
definition in one dataset obtained a perfect score for PPV. Overall, 31
of the 33 studies reported PPVs ≥ 70\%. They also found, however, that
of the 33 studies only 11 reported sensitivity and 5 reported
specificity. PPV is a function of sensitivity, specificity, and
prevalence. Datasets with different values for prevalence will produce
different values for PPV with sensitivity and specificity held constant.
Without sensitivity and specificity, correcting for bias due to
imperfect cohort definitions is not possible. Additionally, the
misclassification of the health condition may be differential, meaning
the cohort definition performs differently on one group of persons
relative to the comparison group, or non-differentially, when the cohort
definition performs similarly on both comparison groups. Prior cohort
definition validation studies have not tested for potential differential
misclassification, even though it can lead to strong bias in effect
estimates.

Once the performance metrics have been established for the cohort
definition, these may be used to adjust the results for studies using
these definitions. In theory, adjusting study results for these
measurement error estimates has been well established. In practice,
though, because of the difficulty in obtaining the performance
characteristics, these adjustments are rarely considered. The methods
used to determine the gold standard are described in the remainder of
this section.

\section{Source Record Verification}\label{source-record-verification}

A common method used to validate cohort definitions has been clinical
adjudication through source record verification: a thorough examination
of a person's records by one or more domain experts with sufficient
knowledge to competently classify the clinical condition or
characteristic of interest. Chart review generally follows the following
steps:

\begin{enumerate}
\def\labelenumi{\arabic{enumi}.}
\item
  Obtain permission from local institutional review board (IRB) and/or
  persons as needed to conduct study including chart review.
\item
  Generate cohort using cohort definition to be evaluated. Sample a
  subset of the persons to manually review if there are insufficient
  resources to adjudicate the entire cohort.
\item
  Identify one or more persons with sufficient clinical expertise to
  review person records.
\item
  Determine guidelines for adjudicating whether a person is positive or
  negative for the desired clinical condition or characteristic.
\item
  Clinical experts review and adjudicate all available data for the
  persons within the sample to classify each person as to whether they
  belong to the phenotype or not.
\item
  Tabulate persons according to the cohort definition classification and
  clinical adjudication classification into a confusion matrix, and
  calculate the performance characteristics possible from the data
  collected.
\end{enumerate}

Results from a chart review are typically limited to the evaluation of
one performance characteristic, positive predictive value (PPV). This is
because the cohort definition under evaluation only generates persons
that are believed to have the desired condition or characteristics.
Therefore, each person in the sample of the cohort is classified as
either a true positive or false positive based on the clinical
adjudication. Without knowledge of all persons in the phenotype in the
entire population (including those not identified by the cohort
definition), it is not possible to identify the false negatives, and
thereby fill in the remainder of the confusion matrix to generate the
remaining performance characteristics. Potential methods of identifying
all persons in the phenotype across the population include chart review
of the entire database, which is generally not feasible unless the
overall population is small, or the utilization of comprehensive
clinical registries in which all true cases have already been flagged
and adjudicated, such as tumor registries (see example below).
Alternatively, one can sample persons who do not qualify for the cohort
definition to produce a subset of predicted negatives, and then
repeating steps 3-6 of the chart review above to check whether these
patients are truly lacking the clinical condition or characteristic of
interest can identify true negatives or false negatives. This would
allow the estimation of negative predictive value (NPV), and if an
appropriate estimate of the phenotype prevalence is available, then
sensitivity and specificity can be estimated.

There are a number of limitations to clinical adjudication through
source record verification. As alluded to earlier, chart review can be a
very time-consuming and resource-intensive process, even just for the
evaluation of a single metric such as PPV. This limitation significantly
impedes the practicality of evaluating an entire population to fill out
a complete confusion matrix. In addition, multiple steps in the above
process have the potential to bias the results of the study. For
example, if records are not equally accessible in the EHR, if there is
no EHR, or if individual patient consent is required, then the subset
under evaluation may not be truly random and could introduce sampling or
selection bias. In addition, manual adjudication is susceptible to human
error or misclassification and thereby may not represent a perfectly
accurate metric. There can often be disagreement between clinical
adjudicators due to the data in the person's record being vague,
subjective, or of low quality. In many studies, the process involves a
majority-rules decision for consensus which yields a binary
classification for persons that does not reflect the inter-rater
discordance.

\subsection{Example of Source Record
Verification}\label{example-of-source-record-verification}

An example of the process to conduct a cohort definition validation
using chart review is provided from a study by the Columbia University
Irving Medical Center (CUIMC), which validated a cohort definition for
multiple cancers as part of a feasibility study for the National Cancer
Institute (NCI). The steps used to conduct the validation for the
example of one of these cancers---prostate cancer---are as follows:

\begin{enumerate}
\def\labelenumi{\arabic{enumi}.}
\tightlist
\item
  Submitted proposal and obtained IRB consent for OHDSI cancer
  phenotyping study.
\item
  Developed a cohort definition for prostate cancer: Using ATHENA and
  ATLAS to explore the vocabulary, we created a cohort definition to
  include all patients with a condition occurrence for Malignant Tumor
  of Prostate (concept ID 4163261), excluding Secondary Neoplasm of
  Prostate (concept ID 4314337) or Non-Hodgkin's Lymphoma of Prostate
  (concept ID 4048666).
\item
  Generated cohort using ATLAS and randomly selected 100 patients for
  manual review, mapping each PERSON\_ID back to patient MRN using
  mapping tables. 100 patients were selected in order to achieve our
  desired level of statistical precision for the performance metric of
  PPV.
\item
  Manually reviewed records in the various EHRs---both inpatient and
  outpatient---in order to determine whether each person in the random
  subset was a true or false positive.
\item
  Manual review and clinical adjudication were performed by one
  physician (although ideally in future more rigorous validation studies
  would be done by a higher number of reviewers to assess for consensus
  and inter-rater reliability).
\item
  Determination of a reference standard was based on clinical
  documentation, pathology reports, labs, medications and procedures as
  documented in the entirety of the available electronic patient record.
\item
  Patients were labeled as 1) prostate cancer 2) no prostate cancer or
  3) unable to determine.
\item
  A conservative estimate of PPV was calculated using the following:
  prostate cancer/ (no prostate cancer + unable to determine).
\item
  Then, using the tumor registry as an additional gold standard to
  identify a reference standard across the entire CUIMC population, we
  counted the number of persons in the tumor registry which were and
  were not accurately identified by the cohort definition, which allowed
  us to estimate sensitivity using these values as true positives and
  false negatives.
\item
  Using the estimated sensitivity, PPV, and prevalence, we could then
  estimate specificity for this cohort definition. As noted previously,
  this process was time-consuming and labor-intensive, as each cohort
  definition had to be individually evaluated through manual chart
  review as well as correlated with the CUIMC tumor registry in order to
  identify all performance metrics. The IRB approval process itself took
  weeks despite an expedited review while obtaining access to the tumor
  registry, and the process of manual chart review itself took a few
  weeks longer.
\end{enumerate}

A review of validation efforts for myocardial infarction (MI) cohort
definitions by Rubbo et al. found that there was significant
heterogeneity in the cohort definitions used in the studies as well as
in the validation methods and the results reported. (2) The authors
concluded that for acute myocardial infarction there is no gold standard
cohort definition available. They noted that the process was both costly
and time-consuming. Due to that limitation, most studies had small
sample sizes in their validation leading to wide variations in the
estimates for the performance characteristics. They also noted that in
the 33 studies, while all the studies reported positive predictive
value, only 11 studies reported sensitivity and only five studies
reported specificity. As mentioned previously, without estimates of
sensitivity and specificity, statistical correction for
misclassification bias cannot be performed.

\section{PheValuator}\label{phevaluator}

The OHDSI community has developed a different approach to constructing a
gold standard by using diagnostic predictive models.(3,4) The general
idea is to emulate the ascertainment of the health outcome similar to
the way clinicians would in a source record validation, but in an
automated way that can be applied at scale. The tool has been developed
as an open-source R package called PheValuator.{[}1{]} PheValuator uses
functions from the Patient Level Prediction package.

The process is as follows:

\begin{enumerate}
\def\labelenumi{\arabic{enumi}.}
\tightlist
\item
  Create an extremely specific (``\textbf{xSpec}'') cohort: Determine a
  set of persons with a very high likelihood of having the outcome of
  interest to be used as noisy positive labels when training a
  diagnostic predictive model.
\item
  Create an extremely sensitive (``\textbf{xSens}'') cohort: Determine a
  set of persons that should include anyone who could possibly have the
  outcome. This cohort will be used to identify its inverse: the set of
  people we are confident do not have the outcome, to be used as noisy
  negative labels when training a diagnostic predictive model.
\item
  Fit a predictive model using the xSpec and xSens cohort: As described
  in Chapter 13, we fit a model using a wide array of patient features
  as predictors and aim to predict whether a person belongs to the xSpec
  cohort (those we believe have the outcome) or the inverse of the xSens
  cohort (those we believe do not have the outcome).
\item
  Apply the fitted model to estimate the probability of the outcome for
  a hold-out set of persons who will be used to evaluate cohort
  definition performance: The set of predictors from the model can be
  applied to a person's data to estimate the predicted probability that
  the person belongs to the phenotype. We use these predictions as a
  \textbf{probabilistic gold standard}.
\item
  Evaluate the performance characteristics of the cohort definitions: We
  compare the predicted probability to the binary classification of a
  cohort definition (the test conditions for the confusion matrix).
  Using the test conditions and the estimates for the true conditions,
  we can fully populate the confusion matrix and estimate the entire set
  of performance characteristics, i.e., sensitivity, specificity, and
  predictive values.
\end{enumerate}

The primary limitation to using this approach is that the estimation of
the probability of a person having the health outcome is limited by the
data in the database. Depending on the database, important information,
such as clinician notes, may not be available.

In diagnostic predictive modeling we create a model that discriminates
between those with the disease and those without the disease. As
described in the Patient-Level Prediction chapter (Chapter 13),
prediction models are developed using a \emph{target cohort} and an
\emph{outcome cohort}. The target cohort includes persons with and
without the health outcome; the outcome cohort identifies those persons
in the target cohort with the health outcome. For the PheValuator
process, we use an extremely specific cohort definition, the ``xSpec''
cohort, to determine the outcome cohort for the prediction model. The
xSpec cohort uses a definition to find those with a very high
probability of having the disease of interest. The xSpec cohort may be
defined as those persons who have multiple condition occurrence records
for the health outcome of interest in a specified period of time. For
example, for a chronic disease such as atrial fibrillation, we may have
persons who have two or more records with the atrial fibrillation
diagnosis code in a 14-day period. For MI, an acute outcome, we may use
two or more occurrences of MI during a single day and include the
requirement of having at least two occurrences from an inpatient
setting. The target cohort for the predictive model is constructed from
the union of persons with a low likelihood of having the health outcome
of interest and those persons in the xSpec cohort. To determine those
persons with a low likelihood of having the health outcome of interest,
we sample from the entire database and exclude persons who have some
evidence suggestive of belonging to the phenotype, typically by removing
persons with any records containing the concepts used to define the
xSpec cohort. There are limitations to this method. It is possible that
these xSpec cohort persons may have different characteristics than
others with the disease. We use LASSO logistic regression to create the
prediction model used to generate the probabilistic gold standard.(5)
This algorithm produces a parsimonious model and typically removes many
of the collinear covariates which may be present across the dataset. In
the current version of the PheValuator software, outcome status (yes/no)
is evaluated based on parameters set in the analysis specification, For
example, for chronic conditions, we may set the time to determine the
characteristics within 365 days of the start of the condition. To add
greater detail to the model, we may break the 365 day observation time
into three separate time windows, such as 0-30 days, 31-90 days, and 91
to 365 days. For acute conditions, we may limit the time to 30 days
after the diagnosis. PheValuator does not evaluate the accuracy of the
cohort start date.

\subsection{Example Validation By
PheValuator}\label{example-validation-by-phevaluator}

We may use PheValuator to assess the complete performance
characteristics for a cohort definition to be used in a study where it
is necessary to determine those persons who have had an acute myocardial
infarction (MI).

The following are the steps for testing cohort definitions for MI using
PheValuator:

\subsubsection{Step 1: Define the xSpec
Cohort}\label{step-1-define-the-xspec-cohort}

Determine those with MI with a high probability: For the cohort entry
event, we required a condition occurrence record with a concept for
myocardial infarction or any of its descendants
(Figure~\ref{fig-quality-clinical-020-cohort-entry}). We required that
each subject in the xSpec cohort have at least 30 days observation time
after the cohort entry event to ensure that data for the model will be
available through the complete observed prediction window.

\begin{figure}[H]

\centering{

\pandocbounded{\includegraphics[keepaspectratio]{quality/images/clinical/fig-quality-clinical-020-cohort-entry.png}}

}

\caption{\label{fig-quality-clinical-020-cohort-entry}Cohort entry event
in ATLAS for an extremely specific cohort definition (xSpec) for
myocardial infarction.}

\end{figure}%

We next added an inclusion criteria requiring either a drug exposure of
anti-thrombotic agent or a second diagnosis code for myocardial
infarction on the same day as the cohort entry event
(Figure~\ref{fig-quality-clinical-030-cohort-inclusion}). These
inclusion criteria increase the specificity of the cohort, i.e.,
increasing the likelihood that the subjects selected by the definition
is a case of MI.

\begin{figure}[H]

\centering{

\pandocbounded{\includegraphics[keepaspectratio]{quality/images/clinical/fig-quality-clinical-030-cohort-inclusion.png}}

}

\caption{\label{fig-quality-clinical-030-cohort-inclusion}Cohort
inclusion event in ATLAS for an extremely specific cohort definition
(xSpec) for myocardial infarction.}

\end{figure}%

Finally, we add an exclusion event where we exclude a set of
differential diagnoses in the period 7 days before and 14 days after the
cohort entry event
(Figure~\ref{fig-quality-clinical-040-cohort-exclusion}). These include
conditions such as myocarditis and anxiety disorder. These exclusion
events help to increase the validity of the diagnosis of myocardial
infarction, further ensuring that the subjects in this cohort have a
high probability of having the condition of interest.

\begin{figure}[H]

\centering{

\pandocbounded{\includegraphics[keepaspectratio]{quality/images/clinical/fig-quality-clinical-040-cohort-exclusion.png}}

}

\caption{\label{fig-quality-clinical-040-cohort-exclusion}Cohort
exclusion event in ATLAS for an extremely specific cohort definition
(xSpec) for myocardial infarction.}

\end{figure}%

\subsubsection{Step 2: Define the xSens
Cohort}\label{step-2-define-the-xsens-cohort}

Next, we develop an extremely sensitive cohort (xSens). This cohort may
be defined for MI as those persons with at least one condition
occurrence record containing a myocardial infarction concept at any time
in their medical history.
Figure~\ref{fig-quality-clinical-050-cohort-sensitive} illustrates the
xSens cohort definition for MI in ATLAS.

\begin{figure}[H]

\centering{

\pandocbounded{\includegraphics[keepaspectratio]{quality/images/clinical/fig-quality-clinical-050-cohort-sensitive.png}}

}

\caption{\label{fig-quality-clinical-050-cohort-sensitive}An extremely
sensitive cohort definition (xSens) for myocardial infarction.}

\end{figure}%

The xSens cohort will be used to find all subjects with even a low
probability of having the outcome of interest. When we select a large
random set of subjects and remove those in the xSens cohort, the
subjects remaining will likely all have a very low probability of having
the outcome of interest.

\subsubsection{Step 3: Running the PheValuator
process}\label{step-3-running-the-phevaluator-process}

The PheValuator process involves three parts:

\begin{enumerate}
\def\labelenumi{\arabic{enumi}.}
\item
  Developing the diagnostic predictive model
\item
  Applying the model to a large, random set of subjects in the database,
  the ``evaluation cohort'', to produce a ``probabilistic gold
  standard''
\item
  Using the ``probabilistic gold standard'' to calculate the performance
  characteristics for phenotypes to be used in research studies
\end{enumerate}

These steps can be performed using the \emph{runPheValuatorAnalyses}
function as described in the PheValuator vignette stored in the
PheValuator repository in GitHub
(https://github.com/OHDSI/PheValuator/tree/main).

At the end of the process, PheValuator will produce an output file with
the performance characteristics for the tested phenotypes. The desired
performance characteristics for each may depend on the intended use of
the cohort to address the research question of interest. For certain
questions, a very sensitive algorithm may be required; others may
require a more specific algorithm. The process for determining the
performance characteristics for a cohort definition using PheValuator is
shown in
Figure~\ref{fig-quality-clinical-060-phevaluator-characteristic}.

\begin{figure}[H]

\centering{

\pandocbounded{\includegraphics[keepaspectratio]{quality/images/clinical/fig-quality-clinical-060-phevaluator-characteristic.png}}

}

\caption{\label{fig-quality-clinical-060-phevaluator-characteristic}Determining
the Performance Characteristics of a cohort definition using
PheValuator. p(O) = Probability of outcome; TP = True Positive; FN =
False Negative; TN = True Negative; FP = False Positive.}

\end{figure}%

In part A of
Figure~\ref{fig-quality-clinical-060-phevaluator-characteristic}, we
examined the persons from the cohort definition to be tested and found
those persons from the evaluation cohort (created in the previous step)
who were included in the cohort definition (Person IDs 016, 019, 022,
023, and 025) and those from the evaluation cohort who were excluded
from the cohort definition (Person Ids 017, 018, 020, 021, and 024). For
each of these included/excluded persons, we had previously determined
the probability of the health outcome using the predictive model (p(O)).

We estimated the values for True Positives, True Negatives, False
Positives, and False Negatives as follows (Part B of
Figure~\ref{fig-quality-clinical-060-phevaluator-characteristic}):

\begin{enumerate}
\def\labelenumi{\arabic{enumi}.}
\item
  If the cohort definition included a person from the evaluation cohort,
  i.e., the cohort definition considered the person a ``positive.'' The
  predicted probability for the health outcome indicated the expected
  value of the number of counts contributed by that person to the True
  Positives, and one minus the probability indicated the expected value
  of the number of counts contributed by that person to the False
  Positives for that person. We added all the expected values of counts
  across persons to get the total expected value. For example, PersonId
  016 had a predicted probability of 99\% for the presence of the health
  outcome, 0.99 was added to the True Positives (expected value of
  counts added 0.99) and 1.00--0.99 = 0.01 was added to the False
  Positives (0.01 expected value). Another way to think of this is that
  the cohort definition that selected this person got it 99\% right and
  1\% wrong. This was repeated for all the persons from the evaluation
  cohort included in the cohort definition (i.e., PersonIds 019, 022,
  023, and 025).
\item
  Similarly, if the cohort definition did not include a person from the
  evaluation cohort, i.e.~the cohort definition considered the person a
  ``negative,'' one minus the predicted probability for the phenotype
  for that person was the expected value of counts contributed to True
  Negatives and was added to it, and, in parallel, the predicted
  probability for the phenotype was the expected value of counts
  contributed to the False Negatives and was added to it. For example,
  PersonId 017 had a predicted probability of 1\% for the presence of
  the health outcome (and, correspondingly, 99\% for the absence of the
  health outcome) and 1.00 -- 0.01 = 0.99 was added to the True
  Negatives and 0.01 was added to the False Negatives. This was repeated
  for all the persons from the evaluation cohort not included in the
  cohort definition (i.e., PersonIds 018, 020, 021, and 024).
\end{enumerate}

After adding these values over the full set of persons in the evaluation
cohort, we filled the four cells of the confusion matrix with the
expected values of counts for each cell, and we were able to create
point estimates for the tested cohort's performance characteristics,
i.e., sensitivity, specificity, and positive and negative (NPV)
predictive value (Figure 1C). In the example, the sensitivity,
specificity, PPV, and NPV were 0.99, 0.63, 0.42, and 0.99, respectively.
PheValuator also calculated the confidence intervals from these
estimates.

The desired performance characteristics may depend on the intended use
of the cohort to address the research question of interest. For certain
questions, a very sensitive algorithm may be required; others may
require a more specific algorithm. An example of the output from an
analysis for acute myocardial infarction is shown in
Figure~\ref{fig-quality-clinical-070-phevaluator-analysis}.

\begin{figure}[H]

\centering{

\pandocbounded{\includegraphics[keepaspectratio]{quality/images/clinical/fig-quality-clinical-070-phevaluator-analysis.png}}

}

\caption{\label{fig-quality-clinical-070-phevaluator-analysis}Example
output from a PheValuator analysis for acute myocardial infarction.}

\end{figure}%

In this example, we included the results from the xSpec cohort (cohort
ID 11081) as well as the cohort of interest (cohort ID 2072). The
performance characteristics of the xSpec cohort showed high PPV and low
sensitivity compared to the test cohort, ``{[}PL{]} All events of Acute
Myocardial Infarction, inpatient setting with washout period of 365
days''. This is expected as the criteria for the xSpec cohort was very
specific for acute myocardial infarction leading to a high PPV.
PheValuator calculates the F1 score which is the harmonic mean of the
sensitivity and the PPV.

\section{Generalizability of the
Evidence}\label{generalizability-of-the-evidence}

While a cohort can be well-defined and fully evaluated within the
context of a given observational database, the clinical validity is
limited by the extent to which the results are considered generalizable
to the target population of interest. Multiple observational studies on
the same topic can yield different results, which can be caused by not
only by their designs and analytic methods, but also bt their choice of
data source. Madigan et al.~(\hyperref[ref-madigan_2013]{2013})
demonstrated that choice of database affects the result of observational
study. They systematically investigated heterogeneity in the results for
53 drug-outcome pairs and two study designs (cohort studies and
self-controlled case series) across the 10 observational databases. Even
though they held study design constant, substantial heterogeneity in
effect estimates was observed.

Across the OHDSI network, observational databases vary considerably in
the populations they represent (e.g.~pediatric vs.~elderly,
privately-insured employees vs.~publicly-insured unemployed), the care
settings where data are captured (e.g.~inpatient vs.~outpatient, primary
vs.~secondary/specialty care), the data capture processes
(e.g.~administrative claims, EHRs, clinical registries), and the
national and regional health system from which care is based. These
differences can manifest as heterogeneity observed when studying disease
and the effects of medical interventions and can also influence the
confidence we have in the quality of each data source that may
contribute evidence within a network study. While all databases within
the OHDSI network are standardized to the CDM, it is important to
reinforce that standardization does not reduce the true inherent
heterogeneity that is present across populations, but simply provides a
consistent framework to investigate and better understand the
heterogeneity across the network. The OHDSI research network provides
the environment to apply the same analytic process on various databases
across the world, so that researchers can interpret results across
multiple data sources while holding other methodological aspects
constant. OHDSI's collaborative approach to open science in network
research, where researchers across participating data partners work
together alongside those with clinical domain knowledge and
methodologists with analytical expertise, is one way of reaching a
collective level of understanding of the clinical validity of data
across a network that should serve as a foundation for building
confidence in the evidence generated using these data.

\section{Summary}\label{summary-1}

\begin{itemize}
\tightlist
\item
  Clinical validity can be established by understanding the
  characteristics of the underlying data source, evaluating the
  performance characteristics of the cohorts within an analysis, and
  assessing the generalizability of the study to the target population
  of interest.
\item
  A cohort definition can be evaluated on the extent to which persons
  identified in the cohort based on the cohort definition and the
  available observational data accurately reflect the persons who truly
  belong to the phenotype.
\item
  Cohort definition validation requires estimating multiple performance
  characteristics, including sensitivity, specificity, and positive
  predictive value, to fully summarize and enable adjustment for
  measurement error.
\item
  Clinical adjudication through source record verification and
  PheValuator represent two alternative approaches to estimating cohort
  definition validation.
\item
  OHDSI network studies provide a mechanism to examine data source
  heterogeneity and expand the generalizability of findings to improve
  clinical validity of real-world evidence.
\end{itemize}

\section{References}\label{references}

1. Hripcsak G, Albers DJ. High-fidelity phenotyping: richness and
freedom from bias. J Am Med Inform Assoc. 2018 Mar 1;25(3):289--94.

2. Rubbo B, Fitzpatrick NK, Denaxas S, Daskalopoulou M, Yu N, Patel RS,
et al.~Use of electronic health records to ascertain, validate and
phenotype acute myocardial infarction: A systematic review and
recommendations. Int J Cardiol. 2015 May;187:705--11.

3. Swerdel JN, Hripcsak G, Ryan PB. PheValuator: Development and
evaluation of a phenotype algorithm evaluator. J Biomed Inform. 2019
Sep;97:103258.

4. Swerdel JN, Schuemie M, Murray G, Ryan PB. PheValuator 2.0:
Methodological improvements for the PheValuator approach to
semi-automated phenotype algorithm evaluation. J Biomed Inform. 2022
Nov;135:104177.

5. Suchard MA, Simpson SE, Zorych I, Ryan P, Madigan D. Massive
Parallelization of Serial Inference Algorithms for a Complex Generalized
Linear Model. ACM Trans Model Comput Simul. 2013 Jan;23(1):1--17.

{[}1{]} \url{https://github.com/OHDSI/PheValuator}

\chapter{Software Validity}\label{sec-quality-software}

TODO- write abstract

\hfill\break

\textbf{Chapter Leads}: Martijn Schuemie, Anthony Sena

\begin{tcolorbox}[enhanced jigsaw, opacitybacktitle=0.6, coltitle=black, colback=white, bottomrule=.15mm, bottomtitle=1mm, title=\textcolor{quarto-callout-note-color}{\faInfo}\hspace{0.5em}{Note}, titlerule=0mm, rightrule=.15mm, colbacktitle=quarto-callout-note-color!10!white, colframe=quarto-callout-note-color-frame, toptitle=1mm, arc=.35mm, left=2mm, breakable, opacityback=0, leftrule=.75mm, toprule=.15mm]

This is page is currently a stub. The chapter is being written in the
OHDSI Teams directory. When the draft is complete, it will be converted
to markdown and moved to this file.

Author Resources (requires an OHDSI Teams account):

\begin{itemize}
\tightlist
\item
  \href{https://ohdsiorg.sharepoint.com/sites/Workgroup-EducationWorkingGroup/Shared\%20Documents/Forms/AllItems.aspx?id=\%2Fsites\%2FWorkgroup\%2DEducationWorkingGroup\%2FShared\%20Documents\%2FSection\%20IV\%20\%2D\%20Quality\%2FChapter\%2018\%20\%2D\%20Software\%20Validity&viewid=05fec2cc\%2Dec8a\%2D4d04\%2Db565\%2Dcf1289b96f67}{Chapter
  Directory}
\item
  \href{https://ohdsiorg.sharepoint.com/:x:/r/sites/Workgroup-EducationWorkingGroup/_layouts/15/Doc2.aspx?action=edit&sourcedoc=\%7B1fa31e39-1c5f-4918-b878-609ebd9810b3\%7D&wdOrigin=TEAMS-WEB.teamsSdk_ns.rwc&wdExp=TEAMS-TREATMENT&wdhostclicktime=1748104477731&web=1}{Book
  Layout}
\item
  \href{https://ohdsiorg.sharepoint.com/sites/Workgroup-EducationWorkingGroup/Shared\%20Documents/Forms/AllItems.aspx?viewid=05fec2cc\%2Dec8a\%2D4d04\%2Db565\%2Dcf1289b96f67}{Education
  Working Group SharePoint Drive}
\end{itemize}

Public Resources:

\begin{itemize}
\tightlist
\item
  \href{https://ohdsi.github.io/TheBookOfOhdsi/}{Book of OHDSI, Edition
  1}
\item
  \href{https://github.com/OHDSI/TheBookOfOhdsi}{Source Code} for Book
  of OHDSI, Edition 1
\item
  \href{https://ohdsi.org/}{OHDSI Home Page}
\end{itemize}

\end{tcolorbox}

\chapter{Diagnostics}\label{sec-quality-diagnostics}

TODO- write abstract

\hfill\break

\textbf{Chapter Leads}: Martijn Schuemie, Mitch Conover

\begin{tcolorbox}[enhanced jigsaw, opacitybacktitle=0.6, coltitle=black, colback=white, bottomrule=.15mm, bottomtitle=1mm, title=\textcolor{quarto-callout-note-color}{\faInfo}\hspace{0.5em}{Note}, titlerule=0mm, rightrule=.15mm, colbacktitle=quarto-callout-note-color!10!white, colframe=quarto-callout-note-color-frame, toptitle=1mm, arc=.35mm, left=2mm, breakable, opacityback=0, leftrule=.75mm, toprule=.15mm]

This is page is currently a stub. The chapter is being written in the
OHDSI Teams directory. When the draft is complete, it will be converted
to markdown and moved to this file.

Author Resources (requires an OHDSI Teams account):

\begin{itemize}
\tightlist
\item
  \href{https://ohdsiorg.sharepoint.com/sites/Workgroup-EducationWorkingGroup/Shared\%20Documents/Forms/AllItems.aspx?id=\%2Fsites\%2FWorkgroup\%2DEducationWorkingGroup\%2FShared\%20Documents\%2FSection\%20IV\%20\%2D\%20Quality\%2FChapter\%2019\%20\%2D\%20Diagnostics&viewid=05fec2cc\%2Dec8a\%2D4d04\%2Db565\%2Dcf1289b96f67}{Chapter
  Directory}
\item
  \href{https://ohdsiorg.sharepoint.com/:x:/r/sites/Workgroup-EducationWorkingGroup/_layouts/15/Doc2.aspx?action=edit&sourcedoc=\%7B1fa31e39-1c5f-4918-b878-609ebd9810b3\%7D&wdOrigin=TEAMS-WEB.teamsSdk_ns.rwc&wdExp=TEAMS-TREATMENT&wdhostclicktime=1748104477731&web=1}{Book
  Layout}
\item
  \href{https://ohdsiorg.sharepoint.com/sites/Workgroup-EducationWorkingGroup/Shared\%20Documents/Forms/AllItems.aspx?viewid=05fec2cc\%2Dec8a\%2D4d04\%2Db565\%2Dcf1289b96f67}{Education
  Working Group SharePoint Drive}
\end{itemize}

Public Resources:

\begin{itemize}
\tightlist
\item
  \href{https://ohdsi.github.io/TheBookOfOhdsi/}{Book of OHDSI, Edition
  1}
\item
  \href{https://github.com/OHDSI/TheBookOfOhdsi}{Source Code} for Book
  of OHDSI, Edition 1
\item
  \href{https://ohdsi.org/}{OHDSI Home Page}
\end{itemize}

\end{tcolorbox}

\part{OHDSI Research}

\chapter{Study Steps}\label{sec-research-steps}

TODO- write abstract

\hfill\break

\textbf{Chapter Leads}: Anthony Sena

\begin{tcolorbox}[enhanced jigsaw, opacitybacktitle=0.6, coltitle=black, colback=white, bottomrule=.15mm, bottomtitle=1mm, title=\textcolor{quarto-callout-note-color}{\faInfo}\hspace{0.5em}{Note}, titlerule=0mm, rightrule=.15mm, colbacktitle=quarto-callout-note-color!10!white, colframe=quarto-callout-note-color-frame, toptitle=1mm, arc=.35mm, left=2mm, breakable, opacityback=0, leftrule=.75mm, toprule=.15mm]

This is page is currently a stub. The chapter is being written in the
OHDSI Teams directory. When the draft is complete, it will be converted
to markdown and moved to this file.

Author Resources (requires an OHDSI Teams account):

\begin{itemize}
\tightlist
\item
  \href{https://ohdsiorg.sharepoint.com/sites/Workgroup-EducationWorkingGroup/Shared\%20Documents/Forms/AllItems.aspx?id=\%2Fsites\%2FWorkgroup\%2DEducationWorkingGroup\%2FShared\%20Documents\%2FSection\%20V\%20\%2D\%20Studies\%2FChapter\%2020\%20Study\%20Steps&viewid=05fec2cc\%2Dec8a\%2D4d04\%2Db565\%2Dcf1289b96f67}{Chapter
  Directory}
\item
  \href{https://ohdsiorg.sharepoint.com/:x:/r/sites/Workgroup-EducationWorkingGroup/_layouts/15/Doc2.aspx?action=edit&sourcedoc=\%7B1fa31e39-1c5f-4918-b878-609ebd9810b3\%7D&wdOrigin=TEAMS-WEB.teamsSdk_ns.rwc&wdExp=TEAMS-TREATMENT&wdhostclicktime=1748104477731&web=1}{Book
  Layout}
\item
  \href{https://ohdsiorg.sharepoint.com/sites/Workgroup-EducationWorkingGroup/Shared\%20Documents/Forms/AllItems.aspx?viewid=05fec2cc\%2Dec8a\%2D4d04\%2Db565\%2Dcf1289b96f67}{Education
  Working Group SharePoint Drive}
\end{itemize}

Public Resources:

\begin{itemize}
\tightlist
\item
  \href{https://ohdsi.github.io/TheBookOfOhdsi/}{Book of OHDSI, Edition
  1}
\item
  \href{https://github.com/OHDSI/TheBookOfOhdsi}{Source Code} for Book
  of OHDSI, Edition 1
\item
  \href{https://ohdsi.org/}{OHDSI Home Page}
\end{itemize}

\end{tcolorbox}

\chapter{OHDSI Network Research}\label{sec-research-network}

TODO- write abstract

\hfill\break

\textbf{Chapter Leads}: Kristin Kostka

\begin{tcolorbox}[enhanced jigsaw, opacitybacktitle=0.6, coltitle=black, colback=white, bottomrule=.15mm, bottomtitle=1mm, title=\textcolor{quarto-callout-note-color}{\faInfo}\hspace{0.5em}{Note}, titlerule=0mm, rightrule=.15mm, colbacktitle=quarto-callout-note-color!10!white, colframe=quarto-callout-note-color-frame, toptitle=1mm, arc=.35mm, left=2mm, breakable, opacityback=0, leftrule=.75mm, toprule=.15mm]

This is page is currently a stub. The chapter is being written in the
OHDSI Teams directory. When the draft is complete, it will be converted
to markdown and moved to this file.

Author Resources (requires an OHDSI Teams account):

\begin{itemize}
\tightlist
\item
  \href{https://ohdsiorg.sharepoint.com/sites/Workgroup-EducationWorkingGroup/Shared\%20Documents/Forms/AllItems.aspx?id=\%2Fsites\%2FWorkgroup\%2DEducationWorkingGroup\%2FShared\%20Documents\%2FSection\%20V\%20\%2D\%20Studies\%2FChapter\%2021\%20OHDSI\%20Network\%20Research&viewid=05fec2cc\%2Dec8a\%2D4d04\%2Db565\%2Dcf1289b96f67}{Chapter
  Directory}
\item
  \href{https://ohdsiorg.sharepoint.com/:x:/r/sites/Workgroup-EducationWorkingGroup/_layouts/15/Doc2.aspx?action=edit&sourcedoc=\%7B1fa31e39-1c5f-4918-b878-609ebd9810b3\%7D&wdOrigin=TEAMS-WEB.teamsSdk_ns.rwc&wdExp=TEAMS-TREATMENT&wdhostclicktime=1748104477731&web=1}{Book
  Layout}
\item
  \href{https://ohdsiorg.sharepoint.com/sites/Workgroup-EducationWorkingGroup/Shared\%20Documents/Forms/AllItems.aspx?viewid=05fec2cc\%2Dec8a\%2D4d04\%2Db565\%2Dcf1289b96f67}{Education
  Working Group SharePoint Drive}
\end{itemize}

Public Resources:

\begin{itemize}
\tightlist
\item
  \href{https://ohdsi.github.io/TheBookOfOhdsi/}{Book of OHDSI, Edition
  1}
\item
  \href{https://github.com/OHDSI/TheBookOfOhdsi}{Source Code} for Book
  of OHDSI, Edition 1
\item
  \href{https://ohdsi.org/}{OHDSI Home Page}
\end{itemize}

\end{tcolorbox}

\chapter{Engagement with Networks}\label{sec-research-engagement}

TODO- write abstract

\hfill\break

\textbf{Chapter Leads}: Clair Blacketer

\begin{tcolorbox}[enhanced jigsaw, opacitybacktitle=0.6, coltitle=black, colback=white, bottomrule=.15mm, bottomtitle=1mm, title=\textcolor{quarto-callout-note-color}{\faInfo}\hspace{0.5em}{Note}, titlerule=0mm, rightrule=.15mm, colbacktitle=quarto-callout-note-color!10!white, colframe=quarto-callout-note-color-frame, toptitle=1mm, arc=.35mm, left=2mm, breakable, opacityback=0, leftrule=.75mm, toprule=.15mm]

This is page is currently a stub. The chapter is being written in the
OHDSI Teams directory. When the draft is complete, it will be converted
to markdown and moved to this file.

Author Resources (requires an OHDSI Teams account):

\begin{itemize}
\tightlist
\item
  \href{https://ohdsiorg.sharepoint.com/sites/Workgroup-EducationWorkingGroup/Shared\%20Documents/Forms/AllItems.aspx?id=\%2Fsites\%2FWorkgroup\%2DEducationWorkingGroup\%2FShared\%20Documents\%2FSection\%20V\%20\%2D\%20Studies\%2FChapter\%2022\%20Networks&viewid=05fec2cc\%2Dec8a\%2D4d04\%2Db565\%2Dcf1289b96f67}{Chapter
  Directory}
\item
  \href{https://ohdsiorg.sharepoint.com/:x:/r/sites/Workgroup-EducationWorkingGroup/_layouts/15/Doc2.aspx?action=edit&sourcedoc=\%7B1fa31e39-1c5f-4918-b878-609ebd9810b3\%7D&wdOrigin=TEAMS-WEB.teamsSdk_ns.rwc&wdExp=TEAMS-TREATMENT&wdhostclicktime=1748104477731&web=1}{Book
  Layout}
\item
  \href{https://ohdsiorg.sharepoint.com/sites/Workgroup-EducationWorkingGroup/Shared\%20Documents/Forms/AllItems.aspx?viewid=05fec2cc\%2Dec8a\%2D4d04\%2Db565\%2Dcf1289b96f67}{Education
  Working Group SharePoint Drive}
\end{itemize}

Public Resources:

\begin{itemize}
\tightlist
\item
  \href{https://ohdsi.github.io/TheBookOfOhdsi/}{Book of OHDSI, Edition
  1}
\item
  \href{https://github.com/OHDSI/TheBookOfOhdsi}{Source Code} for Book
  of OHDSI, Edition 1
\item
  \href{https://ohdsi.org/}{OHDSI Home Page}
\end{itemize}

\end{tcolorbox}

\part{OHDSI in Action}

\chapter{Study Steps}\label{sec-action-therapeutic}

TODO- write abstract

\hfill\break

\textbf{Chapter Leads}: Julie Green

\begin{tcolorbox}[enhanced jigsaw, opacitybacktitle=0.6, coltitle=black, colback=white, bottomrule=.15mm, bottomtitle=1mm, title=\textcolor{quarto-callout-note-color}{\faInfo}\hspace{0.5em}{Note}, titlerule=0mm, rightrule=.15mm, colbacktitle=quarto-callout-note-color!10!white, colframe=quarto-callout-note-color-frame, toptitle=1mm, arc=.35mm, left=2mm, breakable, opacityback=0, leftrule=.75mm, toprule=.15mm]

This is page is currently a stub. The chapter is being written in the
OHDSI Teams directory. When the draft is complete, it will be converted
to markdown and moved to this file.

Author Resources (requires an OHDSI Teams account):

\begin{itemize}
\tightlist
\item
  \href{https://ohdsiorg.sharepoint.com/sites/Workgroup-EducationWorkingGroup/Shared\%20Documents/Forms/AllItems.aspx?id=\%2Fsites\%2FWorkgroup\%2DEducationWorkingGroup\%2FShared\%20Documents\%2FSection\%20VI\%20\%2D\%20Special\%20Therapeutic\%20Areas\%2FChapter\%2023\%20Special\%20Therapeutic\%20Areas&viewid=05fec2cc\%2Dec8a\%2D4d04\%2Db565\%2Dcf1289b96f67}{Chapter
  Directory}
\item
  \href{https://ohdsiorg.sharepoint.com/:x:/r/sites/Workgroup-EducationWorkingGroup/_layouts/15/Doc2.aspx?action=edit&sourcedoc=\%7B1fa31e39-1c5f-4918-b878-609ebd9810b3\%7D&wdOrigin=TEAMS-WEB.teamsSdk_ns.rwc&wdExp=TEAMS-TREATMENT&wdhostclicktime=1748104477731&web=1}{Book
  Layout}
\item
  \href{https://ohdsiorg.sharepoint.com/sites/Workgroup-EducationWorkingGroup/Shared\%20Documents/Forms/AllItems.aspx?viewid=05fec2cc\%2Dec8a\%2D4d04\%2Db565\%2Dcf1289b96f67}{Education
  Working Group SharePoint Drive}
\end{itemize}

Public Resources:

\begin{itemize}
\tightlist
\item
  \href{https://ohdsi.github.io/TheBookOfOhdsi/}{Book of OHDSI, Edition
  1}
\item
  \href{https://github.com/OHDSI/TheBookOfOhdsi}{Source Code} for Book
  of OHDSI, Edition 1
\item
  \href{https://ohdsi.org/}{OHDSI Home Page}
\end{itemize}

\end{tcolorbox}

\chapter{Generative AI}\label{sec-action-genai}

TODO- write abstract

\hfill\break

\textbf{Chapter Leads}: tbc

\begin{tcolorbox}[enhanced jigsaw, opacitybacktitle=0.6, coltitle=black, colback=white, bottomrule=.15mm, bottomtitle=1mm, title=\textcolor{quarto-callout-note-color}{\faInfo}\hspace{0.5em}{Note}, titlerule=0mm, rightrule=.15mm, colbacktitle=quarto-callout-note-color!10!white, colframe=quarto-callout-note-color-frame, toptitle=1mm, arc=.35mm, left=2mm, breakable, opacityback=0, leftrule=.75mm, toprule=.15mm]

This is page is currently a stub. The chapter is being written in the
OHDSI Teams directory. When the draft is complete, it will be converted
to markdown and moved to this file.

Author Resources (requires an OHDSI Teams account):

\begin{itemize}
\tightlist
\item
  \href{https://ohdsiorg.sharepoint.com/sites/Workgroup-EducationWorkingGroup/Shared\%20Documents/Forms/AllItems.aspx?id=\%2Fsites\%2FWorkgroup\%2DEducationWorkingGroup\%2FShared\%20Documents\%2FSection\%20VI\%20\%2D\%20Special\%20Therapeutic\%20Areas\%2FChapter\%2024\%20Generative\%20AI&viewid=05fec2cc\%2Dec8a\%2D4d04\%2Db565\%2Dcf1289b96f67}{Chapter
  Directory}
\item
  \href{https://ohdsiorg.sharepoint.com/:x:/r/sites/Workgroup-EducationWorkingGroup/_layouts/15/Doc2.aspx?action=edit&sourcedoc=\%7B1fa31e39-1c5f-4918-b878-609ebd9810b3\%7D&wdOrigin=TEAMS-WEB.teamsSdk_ns.rwc&wdExp=TEAMS-TREATMENT&wdhostclicktime=1748104477731&web=1}{Book
  Layout}
\item
  \href{https://ohdsiorg.sharepoint.com/sites/Workgroup-EducationWorkingGroup/Shared\%20Documents/Forms/AllItems.aspx?viewid=05fec2cc\%2Dec8a\%2D4d04\%2Db565\%2Dcf1289b96f67}{Education
  Working Group SharePoint Drive}
\end{itemize}

Public Resources:

\begin{itemize}
\tightlist
\item
  \href{https://ohdsi.github.io/TheBookOfOhdsi/}{Book of OHDSI, Edition
  1}
\item
  \href{https://github.com/OHDSI/TheBookOfOhdsi}{Source Code} for Book
  of OHDSI, Edition 1
\item
  \href{https://ohdsi.org/}{OHDSI Home Page}
\end{itemize}

\end{tcolorbox}

\part{Back Matter}

\chapter*{Glossary}\label{sec-appendix-glossary}
\addcontentsline{toc}{chapter}{Glossary}

\markboth{Glossary}{Glossary}

TODO- write abstract

\hfill\break

\textbf{Chapter Leads}: tbc

\begin{tcolorbox}[enhanced jigsaw, opacitybacktitle=0.6, coltitle=black, colback=white, bottomrule=.15mm, bottomtitle=1mm, title=\textcolor{quarto-callout-note-color}{\faInfo}\hspace{0.5em}{Note}, titlerule=0mm, rightrule=.15mm, colbacktitle=quarto-callout-note-color!10!white, colframe=quarto-callout-note-color-frame, toptitle=1mm, arc=.35mm, left=2mm, breakable, opacityback=0, leftrule=.75mm, toprule=.15mm]

This is page is currently a stub. The chapter is being written in the
OHDSI Teams directory. When the draft is complete, it will be converted
to markdown and moved to this file.

Author Resources (requires an OHDSI Teams account):

\begin{itemize}
\tightlist
\item
  Chapter Directory ?
\item
  \href{https://ohdsiorg.sharepoint.com/:x:/r/sites/Workgroup-EducationWorkingGroup/_layouts/15/Doc2.aspx?action=edit&sourcedoc=\%7B1fa31e39-1c5f-4918-b878-609ebd9810b3\%7D&wdOrigin=TEAMS-WEB.teamsSdk_ns.rwc&wdExp=TEAMS-TREATMENT&wdhostclicktime=1748104477731&web=1}{Book
  Layout}
\item
  \href{https://ohdsiorg.sharepoint.com/sites/Workgroup-EducationWorkingGroup/Shared\%20Documents/Forms/AllItems.aspx?viewid=05fec2cc\%2Dec8a\%2D4d04\%2Db565\%2Dcf1289b96f67}{Education
  Working Group SharePoint Drive}
\end{itemize}

Public Resources:

\begin{itemize}
\tightlist
\item
  \href{https://ohdsi.github.io/TheBookOfOhdsi/}{Book of OHDSI, Edition
  1}
\item
  \href{https://github.com/OHDSI/TheBookOfOhdsi}{Source Code} for Book
  of OHDSI, Edition 1
\item
  \href{https://ohdsi.org/}{OHDSI Home Page}
\end{itemize}

\end{tcolorbox}

\chapter*{Protocol Template}\label{sec-appendix-protocol}
\addcontentsline{toc}{chapter}{Protocol Template}

\markboth{Protocol Template}{Protocol Template}

TODO- write abstract

\hfill\break

\textbf{Chapter Leads}: tbc

\begin{tcolorbox}[enhanced jigsaw, opacitybacktitle=0.6, coltitle=black, colback=white, bottomrule=.15mm, bottomtitle=1mm, title=\textcolor{quarto-callout-note-color}{\faInfo}\hspace{0.5em}{Note}, titlerule=0mm, rightrule=.15mm, colbacktitle=quarto-callout-note-color!10!white, colframe=quarto-callout-note-color-frame, toptitle=1mm, arc=.35mm, left=2mm, breakable, opacityback=0, leftrule=.75mm, toprule=.15mm]

This is page is currently a stub. The chapter is being written in the
OHDSI Teams directory. When the draft is complete, it will be converted
to markdown and moved to this file.

Author Resources (requires an OHDSI Teams account):

\begin{itemize}
\tightlist
\item
  Chapter Directory ?
\item
  \href{https://ohdsiorg.sharepoint.com/:x:/r/sites/Workgroup-EducationWorkingGroup/_layouts/15/Doc2.aspx?action=edit&sourcedoc=\%7B1fa31e39-1c5f-4918-b878-609ebd9810b3\%7D&wdOrigin=TEAMS-WEB.teamsSdk_ns.rwc&wdExp=TEAMS-TREATMENT&wdhostclicktime=1748104477731&web=1}{Book
  Layout}
\item
  \href{https://ohdsiorg.sharepoint.com/sites/Workgroup-EducationWorkingGroup/Shared\%20Documents/Forms/AllItems.aspx?viewid=05fec2cc\%2Dec8a\%2D4d04\%2Db565\%2Dcf1289b96f67}{Education
  Working Group SharePoint Drive}
\end{itemize}

Public Resources:

\begin{itemize}
\tightlist
\item
  \href{https://ohdsi.github.io/TheBookOfOhdsi/}{Book of OHDSI, Edition
  1}
\item
  \href{https://github.com/OHDSI/TheBookOfOhdsi}{Source Code} for Book
  of OHDSI, Edition 1
\item
  \href{https://ohdsi.org/}{OHDSI Home Page}
\end{itemize}

\end{tcolorbox}


\backmatter


\end{document}
